\appendixsection{一个连通性性质}

引理 1. 设 X 为连通拓扑空间，\(\Omega\) 为 X 的稠密开子集。若对任意 \(x \in X\) ，存在 x 的邻域 V 使得 \(V \cap \Omega\) 连通，则 \(\Omega\) 连通。

事实上，设 \(\Omega_0\) 是 \(\Omega\) 的非空既开又闭的子集。取 \(x \in X\) ，并设 \(V\) 是 \(x\) 的邻域，使得 \(V \cap \Omega\) 连通。若 \(x \in \overline{\Omega}_0\) ，

\[
(\mathrm{V} \cap \Omega) \cap \Omega_ {0} = \mathrm{V} \cap \Omega_ {0} \neq \varnothing ,
\]

于是 \(V \cap \Omega \subset \Omega_{0}\) 。这样，由于 \(\Omega\) 在 X 中稠密，\(\overline{\Omega}_{0}\) 是 x 的邻域。因而 \(\overline{\Omega}_{0}\) 非空、既开又闭；又因 X 连通，故 \(\overline{\Omega}_{0} = X\) 。由于 \(\Omega_{0}\) 在 \(\Omega\) 中闭，由此推出 \(\Omega_{0} = \Omega \cap \overline{\Omega}_{0} = \Omega\) ，这就证明了 \(\Omega\) 连通。

引理 2. 设 U 为 \(\mathbf{C}^n\) 中的开球，\(f: \mathrm{U} \to \mathbf{C}\) 为不恒为零的全纯函数。设 A 为 U 的子集，使 \(f = 0\) 在 A 上成立。则 U-A 在 U 中稠密且连通。

U-A 的稠密性由《微分与解析流形（结果）》3.2.5 得出。先设 \(n = 1\) 。若 \(a \in \mathrm{A}\) ，则 \(f\) 在 \(a\) 处的幂级数展开（《微分与解析流形（结果）》3.2.1）不恒为 0，由此推出，在 U 中存在邻域 \(\mathrm{V}_a\) 包含 \(a\) ，使 \(f\) 在 \(\mathrm{V}_a - \{a\}\) 上不取零值。于是 \(a\) 在 A 中孤立，这就证明 A 是 U 的离散子集，从而由于 U 在无穷远处可数，A 是可数的。取 \(x, y \in \mathrm{U} - \mathrm{A}\) 。把 \(x\) （相应地 \(y\) ）与 A 中一点连接起来的实仿射直线的并集是贫集（《一般拓扑学》第九章 §5 第 2 节 定义 2）。因此存在 \(z \in \mathrm{U} - \mathrm{A}\) ，使线段 \([x, z]\) 与 \([y, z]\) 都不与 A 相交。于是三点 \(x, y, z\) 属于 U-A 的同一连通分支，这就证明了 \(n = 1\) 情形的引理。转而考虑一般情形。可以假设 A 是 \(f\) 的零点集（《一般拓扑学》第一章 §11 第 1 节 命题 1）。取 \(x, y \in \mathrm{U} - \mathrm{A}\) ，并设 L 为包含 \(x\) 与 \(y\) 的仿射直线。\(f\) 在 \(\mathrm{L} \cap \mathrm{U}\) 上的限制不恒为零，因为 \(x \in \mathrm{L} \cap \mathrm{U}\) 。由已证者，\(x\) 与 \(y\) 属于 \((\mathrm{L} \cap \mathrm{U}) - (\mathrm{L} \cap \mathrm{A})\) 的同一连通分支，从而属于 U-A 的同一连通分支。

引理 3. 设 X 为有限维连通复解析流形，并设 A 为 X 的满足下述条件的子集：

对任意 \(x \in X\) ，存在在 x 处不取零值的解析函数芽 \(f_{x}\) ，使得 A 在 x 处的芽包含于 \(f_{x}\) 的零点集在 x 处的芽之中。

则 X-A 在 X 中稠密且连通。

X-A 的稠密性由《微分与解析流形（结果）》3.2.5 得出。可以假设 A 是闭的（《一般拓扑学》第一章 §11 第 1 节 命题 1）。对任意 \(x \in \mathrm{X}\) ，存在开邻域 \(\mathrm{V}\) 包含 \(x\) ，以及同构 \(c\) ，它把 \(\mathrm{V}\) 映到 \(\mathbf{C}^n\) 中的一个开球，且使得 \(c(\mathrm{A} \cap \mathrm{V})\) 包含于某个在 \(c(\mathrm{V})\) 上不恒为零的全纯函数的零点集中。于是由引理 2，\(\mathrm{V} \cap (\mathrm{X} - \mathrm{A})\) 连通。依据引理 1，这就证明 \(\mathrm{X} - \mathrm{A}\) 连通。

\specialchapter{习题}

假定所有李代数及其上的模均在 \(k\) 上有限维；从 §3 起，假定 \(k\) 的特征为零。

\exercisesection{\S~1}

\begin{enumerate}
\def\labelenumi{\arabic{enumi})}
\item
  假设 \(k\) 的特征为 \(p > 0\) 。设 \(V\) 为向量空间，\(S\) 为有限集。映射 \(r: S \to \operatorname{End}(V)\) 满足拟交换条件 (AC) 当且仅当存在幂 \(q\) （\(p\) 的幂），使得 \([s^q, s'^q] = 0\) 对一切 \(s, s' \in S\) 成立。（利用第一章 §1 习题 19 公式 (1)。）
\item
  假设 \(k\) 是完全域。设 \(V\) 是有限维向量空间，\(u, v \in \operatorname{End}(V)\) 。设 \(u_s, u_n, v_s, v_n\) 是 \(u, v\) 的半单分量与幂零分量。下列条件等价：(i) 存在整数 \(m\) 使 \((\operatorname{ad} u)^m v = 0\) ；(ii) \(u_s\) 与 \(v\) 交换。（为证 (i) \(\Longrightarrow\) (ii)，化归到 \(k\) 代数闭的情形并利用引理 1 (ii)。）
\item
  我们采用第 2 节中的假设。假设 \(k\) 无限且拟交换条件 (AC) 成立。设 \(k'\) 是 \(k\) 的完全扩张。设 \(\lambda : S \to k\) 满足 \(V^{\lambda}(S) \neq 0\) 。置
\end{enumerate}

\[
\mathrm{V} ^ {\prime} = \mathrm{V} \otimes_ {k} k ^ {\prime}, \quad \mathrm{S} ^ {\prime} = \mathrm{S} \otimes_ {k} k ^ {\prime}.
\]

设 \(r' : S' \to \text{End}(V')\) 是由 \(r\) 经纯量扩张得到的线性映射。存在唯一的映射 \(\lambda' : S' \to k'\) 使得 \(V^{\lambda}(S) \otimes_k k' = V'^{\lambda'}(S')\) 。（化归到 \(V = V^{\lambda}(S)\) 的情形。设 \(P\) 是 \(S\) 上的多项式函数，\(q\) 是 \(k\) 的特征指数的一个幂且整除 \(\dim V\) ，使得 \(\lambda^q = P\) 。设 \(P'\) 是 \(S'\) 上扩张 \(P\) 的多项式函数。对每个 \(s' \in S'\) ，存在 \(\lambda'(s') \in k'\) 使得 \(\lambda'(s')^q = P'(s')\) 。证明 \(r'(s')\) 的特征多项式为 \((X - \lambda'(s'))^{\dim V}\) 。）

\begin{enumerate}
\def\labelenumi{\arabic{enumi})}
\setcounter{enumi}{3}
\tightlist
\item
  假设 k 的特征为零。设 \(\mathfrak{g} = \mathfrak{sl}(3, k)\) ，并用 a 表示由特征值为 1、-1、0 的对角矩阵生成的 g 的子代数。证明 a 在 g 中是约化的，a 在 g 中的交换子 m 由迹为零的对角矩阵组成，且 m 在 g 中的交换子等于 m，从而异于 a（参见第 5 节的注）。
\end{enumerate}

¶ 5) 假设 k 无限。设 g 为李代数，V 为 g-模。若 n 是整数 \(\geq 0\) ，用 \(V_{n}\) 表示由 \(v \in V\) 中使 \(x^{n}v = 0\) 对一切 \(x \in g\) 成立者所成的集。

\begin{enumerate}
\def\labelenumi{\alph{enumi})}
\tightlist
\item
  证明：若 \(v \in V_n\) ，\(x, y \in \mathfrak{g}\) ，则
\end{enumerate}

\[
\left(\sum_ {i = 1} ^ {n} x ^ {n - i} y x ^ {i - 1}\right) v = 0.
\]

（利用 \((x + ty)^n v = 0\) 对一切 \(t\in k\) 成立这一事实。）

在此公式中把 \(y\) 换成 \([x, y]\) ，据此\(^3\)推出 \((x^n y - y x^n)v = 0\) ，从而 \(x^n y v = 0\) 。

\begin{enumerate}
\def\labelenumi{\alph{enumi})}
\setcounter{enumi}{1}
\item
  证明 \(V_{n}\) 是 V 的 g-子模（利用 a)）。特别地，\(\mathrm{V}^{0}(\mathfrak{g})=\bigcup_{n}\mathrm{V}_{n}\) 是 V 的 g-子模。
\item
  假设 \(k = \mathbf{R}\) 或 \(\mathbf{C}\) ，用 \(G\) 表示以 \(\mathfrak{g}\) 为李代数的单连通李群；\(\mathfrak{g}\) 在 \(V\) 上的作用定义了 \(G\) 在 \(V\) 上的作用律（第三章 §6 第 1 节）。证明元素 \(v \in V\) 属于 \(V_n\) 当且仅当 \((s - 1)^n v = 0\) 对一切 \(s \in G\) 成立；特别地，\(V^0(\mathfrak{g}) = V^1(G)\) 。
\end{enumerate}

\begin{enumerate}
\def\labelenumi{\arabic{enumi})}
\setcounter{enumi}{5}
\tightlist
\item
  记号取自第一章 §3 习题 12。特别地，g 是李代数，M 是 g-模，而 \(\mathrm{H}^{p}(\mathfrak{g},\mathrm{M})=\mathrm{Z}^{p}(\mathfrak{g},\mathrm{M})/\mathrm{B}^{p}(\mathfrak{g},\mathrm{M})\) 是 g 的以 M 为系数的 p 次上同调空间。
\end{enumerate}

\begin{enumerate}
\def\labelenumi{\alph{enumi})}
\item
  证明 \(\mathrm{B}^{p}(\mathfrak{g},\mathrm{M})\) 与 \(\mathrm{Z}^{p}(\mathfrak{g},\mathrm{M})\) 在 g 的自然表示 \(\theta\) （作用于上链空间 \(\mathrm{C}^{p}(\mathfrak{g},\mathrm{M})\) ）之下稳定。由此得到 g 在 \(\mathrm{H}^{p}(\mathfrak{g},\mathrm{M})\) 上的一个表示。证明此表示是平凡的（利用公式 \(\theta = di + id\) ，见同前）。
\item
  设 \(\overline{k}\) 是 k 的代数闭包。设 \(x \in g\) ，并用 \(x_{M}\) 表示 M 上相应的自同态。设 \(\lambda_{1}, \ldots, \lambda_{n}\) （相应地 \(\mu_{1}, \ldots, \mu_{m}\) ）是 \(\overline{k}\) 中 \(ad_{g}x\) （相应地 \(x_{M}\) ）的特征值，且按重数重复。证明自同态 \(\theta(x)\) 在 \(C^{p}(g, M)\) 上的特征值是诸 \(\mu_{j} - (\lambda_{i_{1}} + \cdots + \lambda_{i_{p}})\) ，其中 \(1 \leq j \leq m\) 且
\end{enumerate}

\[
1 \leq i _ {1} <   i _ {2} <   \dots <   i _ {p} \leq n.
\]

试利用 a) 推出：\(\mathrm{H}^{p}(\mathfrak{g},\mathrm{M})=0\) ，只要诸 \(\mu_{j}-(\lambda_{i_{1}}+\cdots+\lambda_{i_{p}})\) 均不为零。

\begin{enumerate}
\def\labelenumi{\alph{enumi})}
\setcounter{enumi}{2}
\tightlist
\item
  假设表示 \(\mathfrak{g} \to \operatorname{End}(\mathrm{M})\) 是忠实的，且 \(x\) 满足条件：
\end{enumerate}

\[
\mu_ {j _ {1}} + \dots + \mu_ {j _ {p}} \neq \mu_ {k _ {1}} + \dots + \mu_ {k _ {p + 1}}\tag{($S_{p}$)}
\]

对所有 \(j_{1},\ldots,j_{p},k_{1},\ldots,k_{p+1}\in[1,m]\) 。

证明此时有 \(\mathrm{H}^p (\mathfrak{g},\mathrm{M}) = 0\) （注意 \(\lambda_{i}\) ，即 \(\mathrm{ad}_{\mathfrak{g}}x\) 的特征值，形如 \(\mu_j - \mu_k\) ，并应用 \(b\) ）。

\begin{enumerate}
\def\labelenumi{\arabic{enumi})}
\setcounter{enumi}{6}
\tightlist
\item
  设 \(\mathfrak{g}\) 是幂零李代数，V 是 \(\mathfrak{g}\) -模且 \(V^0(\mathfrak{g}) = 0\) 。证明 \(H^p(\mathfrak{g}, V) = 0\) 对一切 \(p \geq 0\) 成立。（化归到 \(V = V^\lambda(\mathfrak{g})\) 且 \(\lambda \neq 0\) 的情形，并选取元素 \(x \in \mathfrak{g}\) 使 \(\lambda(x) \neq 0\) 。应用习题 6 b)，注意诸 \(\lambda_i\) 全为零而诸 \(\mu_j\) 都等于 \(\lambda(x)\) 。）
\end{enumerate}

重新得到命题 9 的推论（取 p = 1）。\(^{4}\)

¶ 8) 假设 k 的特征为 p \textgreater{} 0。设 g 是 k 上的李代数，具有基 \((e_{1},\ldots,e_{n})\) 。用 U 表示 g 的包络代数，C 表示 U 的中心。对 \(i = 1,\ldots,n\) ，选取一个非零 p-多项式 \(f_{i}\) ，其次数为 \(d_{i}\) ，使得 \(f_{i}(\mathrm{ad}e_{i}) = 0\) ；于是 \(f_{i}(e_{i}) \in \mathrm{C}\) ，参见第一章 §7 习题 5。置 \(z_{i} = f_{i}(e_{i})\) 。

\begin{enumerate}
\def\labelenumi{\alph{enumi})}
\item
  证明 \(z_{1},\ldots,z_{n}\) 代数无关。若 \(A=k[z_{1},\ldots,z_{n}]\) ，证明 U 是自由 A-模，以单项式 \(e_{1}^{\alpha_{1}}\ldots e_{n}^{\alpha_{n}}\) （其中 \(0\leq\alpha_{i}\leq d_{i}\) ）为基。（利用 Poincaré-Birkhoff-Witt 定理。）U 在 A 上的秩 {[}U:A{]} 等于 \(d_{1}\ldots d_{n}\) ；它是 p 的幂。由此推出 C 是有限型 A-模，从而是有限型 k-代数且维数为 n（《交换代数》第八章）。
\item
  设 \(\mathrm{K}\) 是 \(\mathrm{A}\) 的分式域，并置
\end{enumerate}

\[
\mathrm{U} _ {(\mathrm{K})} = \mathrm{U} \otimes_ {\mathrm{A}} \mathrm{K}, \quad \mathrm{C} _ {(\mathrm{K})} = \mathrm{C} \otimes_ {\mathrm{A}} \mathrm{K}.
\]

则有 \(\mathrm{U}_{(\mathrm{K})}\supset\mathrm{C}_{(\mathrm{K})}\supset\mathrm{K}\) 。证明 \(\mathrm{U}_{(\mathrm{K})}\) 是一个域，其中心为 \(\mathrm{C}_{(\mathrm{K})}\) ，且它是 U 的（左、右）商域，参见第一章 §2 习题 10。推出 \([\mathrm{U}_{(\mathrm{K})}:\mathrm{C}_{(\mathrm{K})}]\) 具有 \(q^{2}\) 的形式，其中 q 是 p 的幂；我们有 \([\mathrm{C}_{(\mathrm{K})}:\mathrm{K}]=q_{\mathrm{C}}\) ，其中 \(q_{C}\) 是 p 的幂，而 \([U:A]=q_{C}q^{2}\) 。

\begin{enumerate}
\def\labelenumi{\alph{enumi})}
\setcounter{enumi}{2}
\tightlist
\item
  设 d 是 A 的非零元素，\(\Lambda\) 是 \(\mathrm{U}_{(\mathrm{K})}\) 的子环且 \(U \subset \Lambda \subset d^{-1}U\) 。证明 \(\Lambda = U\) 。{[}若 x = b/a（\(a \in A - \{0\}\) ）是 \(\Lambda\) 的元素，对 m 作归纳证明：关系 \(b \in Ua + U_m\) 蕴含 \(b \in Ua + U_{m-1}\) ，其中 \(\{U_m\}\) 是 U 的典则滤链。（为此，利用 gr U 整闭这一事实，并按《交换代数》第五章 §1 第 4 节 命题 15 的方式论证。）对 m = 0，即得 \(b \in Ua\) ，亦即 \(x \in U\) 。{]}
\end{enumerate}

推出 C 是整闭的。

\begin{enumerate}
\def\labelenumi{\alph{enumi})}
\setcounter{enumi}{3}
\tightlist
\item
  假设 k 代数闭。设 \(\rho: \mathfrak{g} \to \mathfrak{gl}(V)\) 是 g 的不可约线性表示，\(\rho_{U}\) 是 U 的相应表示。\(\rho_{U}\) 在 C 上的限制是从 C 到 k 的同态 \(\gamma_{\rho}\) （k 与 V 的位似群等同）；用 \(\alpha_{\rho}\) 表示它在 A 上的限制。证明：对每个从 k-代数 A （相应地 C）到 k 的同态 \(\alpha\) （相应地 \(\gamma\) ），至少存在 g 的一个不可约表示 \(\rho\) 使 \(\alpha_{\rho} = \alpha\) （相应地 \(\gamma_{\rho} = \gamma\) ），且这样的表示（在等价意义下）只有有限多个。证明 \(\dim V \leq q\) ，记号同 b)。\(^{5}\)
\end{enumerate}

¶ 9) 沿用上一习题的记号，并进一步假设 g 是幂零的。

\begin{enumerate}
\def\labelenumi{\alph{enumi})}
\item
  证明可以选取基 \((e_{1},\ldots,e_{n})\) ，使得对任意一对 \((i,j)\) ，\([e_{i},e_{j}]\) 是诸 \(e_{h}\) （ \(h > \sup(i,j)\) ）的线性组合。以下假设诸 \(e_{i}\) 满足此条件。对 \(i=1,\ldots,n\) 选取 p 的幂 \(q(i)\) 使 \(\operatorname{ad}(e_{i})^{q(i)} = 0\) ，并置 \(z_{i} = e_{i}^{q(i)}, A = k[z_{1},\ldots,z_{n}]\) ，参见习题 8. b)。设 \(\rho : g \to gl(V)\) 是 g 的线性表示。假设 \(\rho(e_{i})\) 对 \(i = 1,\ldots,n\) 幂零。证明 \(\rho(x)\) 对一切 \(x \in g\) 幂零。（对 \(n = \dim g\) 作归纳论证，并化归到 \(\rho\) 不可约的情形。证明在此情形 \(\rho(e_{n}) = 0\) ，然后应用归纳假设。）
\item
  设 \(\rho_1: \mathfrak{g} \to \mathfrak{gl}(V_1)\) 与 \(\rho_2: \mathfrak{g} \to \mathfrak{gl}(V_2)\) 是 \(\mathfrak{g}\) 的两个线性表示。假设 \(V_1\) 与 \(V_2\) 均 \(\neq 0\) ，且 \(V_1 = V^{\lambda_1}(\mathfrak{g}), V_2 = V^{\lambda_2}(\mathfrak{g})\) ，其中 \(\lambda_1\) 与 \(\lambda_2\) 是 \(\mathfrak{g}\) 上的两个函数，参见第 3 节。证明：若 \(\lambda_1(e_i) = \lambda_2(e_i)\) 对 \(i = 1, \ldots, n\) 成立，则 \(\lambda_1 = \lambda_2\) ，且存在一个非零 \(\mathfrak{g}\) -同态，从 \(V_1\) 到 \(V_2\) （把 \(b\) 应用于 \(\mathfrak{g}\) -模 \(V = \mathcal{L}(V_1, V_2)\) ，并用 Engel 定理证明 \(V\) 含有非零的 \(\mathfrak{g}\) -不变元）。推出：若此外 \(V_1\) 与 \(V_2\) 还是单模，则它们同构。
\item
  假设 k 代数闭。设 R 是 g 的不可约表示的等价类全体所成的集。若 \(\rho \in R\) ，置
\end{enumerate}

\[
x _ {\rho} = (x _ {\rho} (1), \dots , x _ {\rho} (n)) \in k ^ {n},
\]

其中 \(x_{\rho}(i)\) 是 \(\rho(e_{i})\) 的唯一特征值。证明 \(\rho \mapsto x_{\rho}\) 是从 R 到 \(k^{n}\) 的双射。（单性由 c) 得出，满性由习题 8 d) 得出。）由此推出下列结论：

\begin{enumerate}
\def\labelenumi{(\roman{enumi})}
\item
  对 A 的任一极大理想 \(\mathfrak{m}\) ，U/mU 模其根基的商是矩阵代数。
\item
  \(\mathfrak{g}\) 的任一不可约表示的次数都是 \(p\) 的幂（这由 (i) 及 \([\mathrm{U} / \mathfrak{mU}:k]\) 是 \(p\) 的幂这一事实得出）。
\item
  每个从 A 到 k 的同态都可唯一扩张为从 C 到 k 的同态（利用 C/mC 包含于 U/mU 的中心之中这一事实，而该中心是剩余域为 k 的局部 k-代数）。
\item
  存在整数 N ≥ 0，使 \(x^{p^{N}} \in A\) 对一切 \(x \in C\) 成立（这由 (iii) 得出）。\(^{6}\)
\end{enumerate}

¶ 10) 假设 \(k\) 的特征为 \(p > 0\) 。用 \(\mathfrak{g}\) 表示以 \(\{e_1, e_2, e_3\}\) 为基的李代数，其中 \([e_1, e_2] = e_3\) ，\([e_1, e_3] = [e_2, e_3] = 0\) 。

\begin{enumerate}
\def\labelenumi{\alph{enumi})}
\item
  证明 \(\mathrm{U}\mathfrak{g}\) 的中心是 \(k[e_1^p,e_2^p,e_3]\) 。
\item
  假设 k 代数闭。证明：对一切 \((\lambda_{1},\lambda_{2},\lambda_{3})\in k^{3}\) ，存在（在等价意义下）唯一的 g 的不可约表示 \(\rho\) ，使 \(\lambda_{i}\) 是 \(\rho(e_i)\) 的唯一特征值（ \(i = 1,2,3\) ）；\(\rho\) 的次数为 \(p\) ，若 \(\lambda_3 \neq 0\) ；为 1，若 \(\lambda_3 = 0\) 。（应用习题 8 与 9，或直接论证。）
\end{enumerate}

\begin{enumerate}
\def\labelenumi{\arabic{enumi})}
\setcounter{enumi}{10}
\tightlist
\item
  设 \(\mathfrak{h}\) 为幂零李代数，\(V\) 为不缩为 0 的 \(\mathfrak{h}\) -模，\(\lambda\) 为 \(\mathfrak{h}\) 上的函数且 \(V = V^{\lambda}(\mathfrak{h})\) 。证明下列性质等价：
\end{enumerate}

\begin{enumerate}
\def\labelenumi{(\roman{enumi})}
\item
  \(\lambda\) 是 \(\mathfrak{h}\) 上的线性形式，在 \([\mathfrak{h},\mathfrak{h}]\) 上为零。
\item
  存在 \(\mathrm{V}\) 的一组基，使得由 \(\mathfrak{h}\) 的元素所定义的自同态关于这组基都是三角形的。
\end{enumerate}

（为证 (i) \(\Longrightarrow\) (ii)，对 \(\mathfrak{h}\) -模 \(\mathcal{L}(\mathrm{W}_{\lambda},\mathrm{V})\) 应用 Engel 定理，其中 \(\mathrm{W}\) 是 1 维 \(\mathfrak{h}\) -模，由 \(\lambda\) 所定义。）

当 k 的特征为 0 时，性质 (i) 与 (ii) 成立（命题 9）。

\exercisesection{\S~2}

\begin{enumerate}
\def\labelenumi{\arabic{enumi})}
\item
  迹为零的对角矩阵构成 \(\mathfrak{sl}(n,k)\) 的一个嘉当子代数，但当 \(n = 2\) 且 \(k\) 的特征为 2 时除外。
\item
  设 \(e\) 是元素 \(\begin{pmatrix} 0 & 1 \\ 0 & 0 \end{pmatrix}\) ，它属于 \(\mathfrak{sl}(2, \mathbf{C})\) 。证明 \(\mathbf{C}e\) 是 \(\mathfrak{sl}(2, \mathbf{C})\) 的极大幂零李子代数，但不是 \(\mathfrak{sl}(2, \mathbf{C})\) 的嘉当子代数。
\item
  假设 \(k\) 的特征为 0。设 \(\mathfrak{g}\) 是半单李代数。设 \(E\) 是 \(\mathfrak{g}\) 中交换子代数的集合，其元素均在 \(\mathfrak{g}\) 中半单。则 \(\mathfrak{g}\) 的嘉当子代数正是 \(E\) 的极大元。（利用定理 2 与命题 10。）
\end{enumerate}

特别地，\(\mathfrak{g}\) 的诸嘉当子代数之并等于 \(\mathfrak{g}\) 的半单元素组成的集合。

\begin{enumerate}
\def\labelenumi{\arabic{enumi})}
\setcounter{enumi}{3}
\item
  设 \(\mathfrak{g}\) 是具有基 \((x,y,z)\) 的李代数，使得 \([x,y] = y\) ，\([x,z] = z\) ，\([y,z] = 0\) 。设 \(\mathfrak{a}\) 为理想 \(ky + kz\) （\(\mathfrak{g}\) 中的）。则 \(\mathrm{rk}(\mathfrak{a}) = 2\) 且 \(\mathrm{rk}(\mathfrak{g}) = 1\) 。
\item
  假设 \(k\) 的特征为 0。设 \(\mathfrak{g}\) 是李代数，\(\mathfrak{r}\) 是其根基，\(\mathfrak{h}\) 是 \(\mathfrak{g}\) 的一个嘉当子代数。证明
\end{enumerate}

\[
\mathfrak {r} = [ \mathfrak {g}, \mathfrak {r} ] + (\mathfrak {h} \cap \mathfrak {r}).
\]

（注意 \(\mathfrak{h}\) 在 \(\mathfrak{g}/[\mathfrak{g},\mathfrak{r}]\) 中的像包含 \(\mathfrak{r}/[\mathfrak{g},\mathfrak{r}]\) ，即 \(\mathfrak{g}/[\mathfrak{g},\mathfrak{r}]\) 的中心。）

\begin{enumerate}
\def\labelenumi{\arabic{enumi})}
\setcounter{enumi}{5}
\item
  设 \(\mathfrak{g}\) 是李代数，\(\mathfrak{h}\) 是 \(\mathfrak{g}\) 的幂零子代数。若 \(\mathfrak{g}^0 (\mathfrak{h})\) 幂零，则 \(\mathfrak{g}^0 (\mathfrak{h})\) 是 \(\mathfrak{g}\) 的嘉当子代数。
\item
  设 \(\mathfrak{s}\) 是李代数，\(\mathfrak{a}\) 是 \(\mathfrak{s}\) 的嘉当子代数，V 是 \(\mathfrak{s}\) -模。设 \(\mathfrak{g} = \mathfrak{s} \times V\) 是 \(\mathfrak{s}\) 与 V 的半直积。证明 \(\mathfrak{a} \times V^0(\mathfrak{a})\) 是 \(\mathfrak{g}\) 的嘉当子代数。
\item
  假设 \(k\) 的特征为 \(p > 0\) 。用 \(\mathfrak{s}\) 表示具有基 \(\{x, y\}\) 且满足 \([x, y] = y\) 的李代数。设 \(V\) 是 \(k\) -向量空间，具有基 \(\{e_i\}_{i \in \mathbf{Z}/p\mathbf{Z}}\) 。
\end{enumerate}

\begin{enumerate}
\def\labelenumi{\alph{enumi})}
\item
  证明 V 有唯一的 \(\mathfrak{s}\) -模结构，使得 \(xe_{i} = ie_{i}\) 且 \(ye_{i} = e_{i + 1}\) 对一切 \(i\) 成立。此 \(\mathfrak{s}\) -模是单模。
\item
  设 \(\mathfrak{g} = \mathfrak{s} \times \mathrm{V}\) 是 \(\mathfrak{s}\) 与 \(\mathrm{V}\) 的半直积。证明 \(\mathfrak{g}\) 是秩为 1 的可解代数，其导代数不幂零。
\item
  \(\mathfrak{g}\) 的一个元素是正则元，当且仅当它在 \(\mathfrak{s}\) 上的投影形如 \(ax + by\) ，其中 \(ab \neq 0\) 。
\item
  我们有 \(\mathrm{V}^{0}(x+y)=0\) 且 \(\mathrm{V}^{0}(x)=ke_{0}\) 。推出（参见习题 6）g 有维数为 1 的嘉当子代数（例如由 \(x+y\) 生成的那个），以及维数为 2 的嘉当子代数（例如由 x 与 \(e_{0}\) 生成的那个）。
\end{enumerate}

\begin{enumerate}
\def\labelenumi{\arabic{enumi})}
\setcounter{enumi}{8}
\item
  设 g 是具有基 \((x,y)\) 且满足 \([x,y]=y\) 的李代数。设 k=ky，并用 \(\varphi:g\to g/k\) 表示典则态射。g/k 的元素 0 在 g/k 中正则，但它不是 g 的正则元在 \(\varphi\) 下的像。
\item
  假设 \(k\) 无限。设 \(\mathfrak{g}\) 是李代数。证明下列性质等价：
\end{enumerate}

\begin{enumerate}
\def\labelenumi{(\roman{enumi})}
\item
  \(\operatorname{rk}(\mathfrak{g}) = \dim (\mathfrak{g})\)。
\item
  \(\mathfrak{g}\) 是幂零的。
\item
  g 只有有限多个维数为 rk(g) 的嘉当子代数。
\item
  g 只有一个嘉当子代数。
\end{enumerate}

\begin{enumerate}
\def\labelenumi{\arabic{enumi})}
\setcounter{enumi}{10}
\tightlist
\item
  设 \(\mathfrak{h}\) 是一个 \(\neq 0\) 的交换李代数，\(P\) 是 \(\mathfrak{h}^*\) 的包含 0 的有限子集。证明存在一个李代数 \(\mathfrak{g}\) ，以 \(\mathfrak{h}\) 为嘉当子代数，且使得 \(\mathfrak{h}\) 在 \(\mathfrak{g}\) 中的权集为 \(P\) 。（把 \(\mathfrak{g}\) 构造为 \(\mathfrak{h}\) 与 \(\mathfrak{h}\) -模 \(V\) 的半直积，其中 V 是对应于 \(P - \{0\}\) 中诸元素的一维模的直和，参见习题 7。）
\end{enumerate}

\(x\) 是 \(\mathfrak{h}\) 中的元素，满足 \(\mathfrak{h} = \mathfrak{g}^0(x)\) ，当且仅当 \(x\) 不与 \(\mathrm{P} - \{0\}\) 中任何元素正交。

\begin{enumerate}
\def\labelenumi{\arabic{enumi})}
\setcounter{enumi}{11}
\tightlist
\item
  假设 \(k\) 有限。构造一个李代数 \(\mathfrak{g}\) 的例子，它有一个嘉当子代数 \(\mathfrak{h}\) ，其中不存在元素 \(x\) 使得 \(\mathfrak{h} = \mathfrak{g}^0(x)\) 。（利用前一习题，并取 \(\mathrm{P} = \mathfrak{h}^*\) 。）
\end{enumerate}

¶ 13) 假设 k 有限。用 \(k'\) 表示 k 的一个无限扩张。设 g 是 k 上的李代数。g 的秩记作 \(\operatorname{rk}(\mathfrak{g})\) ，定义为 \(k'\) -李代数 \(g' = g \otimes_{k} k'\) 的秩；g 的元素称为正则的，若它在 \(g'\) 中正则；这些定义不依赖于 \(k'\) 的选取。证明：若

\[
\operatorname{Card} (k) \geq \dim \mathfrak {g} - \operatorname{rk} (\mathfrak {g}),
\]

\(\mathfrak{g}\) 含有正则元（从而也含有嘉当子代数）。

（利用下述结果：若 \(a\) 是 \(k[\mathrm{X}_1, \ldots, \mathrm{X}_n]\) 的非零齐次元素，且 \(\operatorname{Card}(k) \geq \deg(a)\) ，则存在 \(x \in k^n\) 使得 \(a(x) \neq 0\) 。）

\begin{enumerate}
\def\labelenumi{\arabic{enumi})}
\setcounter{enumi}{13}
\tightlist
\item
  假设 \(k\) 的特征为零。设 V 是有限维 \(k\) -向量空间，\(\mathfrak{g}\) 是 \(\mathfrak{gl}(V)\) 的李子代数，\(\mathfrak{h}\) 是 \(\mathfrak{g}\) 的嘉当子代数，\(\mathfrak{n}_{\mathrm{V}}\) 是 \(\mathfrak{g}\) -模 V 的最大幂零理想（第一章 §4 第 3 节 定义 2）。
\end{enumerate}

证明 h 的元素幂零，当且仅当它属于 \(n_{V}\) 。（化归到 g-模 V 半单的情形，并利用定理 2 的推论 3。）

¶ 15) 假设 \(k\) 无限。设 \(\mathfrak{g}\) 是李代数，\(\mathcal{C}_{\infty}\mathfrak{g}\) 是 \(\mathfrak{g}\) 的升中心列之并，\(x\) 是 \(\mathfrak{g}\) 的元素。证明下列性质等价：

\begin{enumerate}
\def\labelenumi{(\roman{enumi})}
\item
  \(x\) 属于 \(\mathfrak{g}\) 的每个嘉当子代数。
\item
  \(x\in \mathfrak{g}^0 (y)\) 对一切 \(y\in \mathfrak{g}\) 成立（即 \(x\in \mathfrak{g}^0 (\mathfrak{g})\) ）。
\item
  \(x \in C_{\infty}g.\)
\end{enumerate}

（蕴涵关系 (iii) \(\Longrightarrow\) (ii) \(\Longrightarrow\) (i) 是直接的。为证 (i) \(\Longrightarrow\) (ii)，注意 (i) 等价于说 \(x\in \mathfrak{g}^0 (y)\) 对每个正则元 \(y\) （\(\mathfrak{g}\) 中的）成立，并利用正则元在 Zariski 拓扑下于 \(\mathfrak{g}\) 中稠密这一事实。为证 (ii) \(\Longrightarrow\) (iii)，注意 \(\mathfrak{n} = \mathfrak{g}^0 (\mathfrak{g})\) 在 \(\mathfrak{g}\) 下稳定（§1 习题 5），并对 \(\mathfrak{g}\) -模 \(\mathfrak{n}\) 应用 Engel 定理；推出 \(\mathfrak{n}\) 含于 \(\mathcal{L}_{\infty}\mathfrak{g}\) 。）

¶ 16) 设 g 是可解复李代数，h 是 g 的嘉当子代数，\(\mathfrak{g} = \bigoplus \mathfrak{g}^{\lambda}(\mathfrak{h})\) 是 g 的相应准素子空间分解，其中 \(\mathfrak{g}^{0}(\mathfrak{h}) = \mathfrak{h}\) 。

\begin{enumerate}
\def\labelenumi{\alph{enumi})}
\item
  证明线性形式在 \(\mathfrak{h}\) 上的限制——这些线性形式即第三章 §9 习题 17 c) 中称为 \(\mathfrak{g}\) 的根者——是 \(\mathfrak{h}\) 在 \(\mathfrak{g}\) 中的权，即诸 \(\lambda\) ，使得 \(\mathfrak{g}^{\lambda}(\mathfrak{h}) \neq 0\) ；推出这样的 \(\lambda\) 在 \(\mathfrak{h} \cap \mathcal{D}\mathfrak{g}\) 上为零。
\item
  设 \((x,y)\mapsto [x,y]'\) 是从 \(\mathfrak{g}\times \mathfrak{g}\) 到 \(\mathfrak{g}\) 的交错双线性映射，具有下列性质：
\end{enumerate}

\begin{enumerate}
\def\labelenumi{(\roman{enumi})}
\item
  若 \(x \in \mathfrak{g}^{\lambda}(\mathfrak{h}), y \in \mathfrak{g}^{\mu}(\mathfrak{h})\) ，其中 \(\lambda \neq 0, \mu \neq 0\) ，则 \([x, y]' = [x, y]\) ；
\item
  若 \(x \in \mathfrak{g}^{0}(\mathfrak{h}), y \in \mathfrak{g}^{\mu}(\mathfrak{h})\) ，则 \([x, y]' = [x, y] - \mu(x)y\) 。
\end{enumerate}

证明这在 \(\mathfrak{g}\) 上给出一个新的李代数结构（利用 \(a\) ）。把它记作 \(\mathfrak{g}'\) 。

\begin{enumerate}
\def\labelenumi{\alph{enumi})}
\setcounter{enumi}{2}
\tightlist
\item
  证明：若 \(x \in \mathfrak{g}^{\lambda}(\mathfrak{h})\) ，则映射 \(ad'x : y \mapsto [x, y]'\) 幂零。推出 \(g'\) 幂零（把第一章 §4 习题 11 应用于 \(ad'x\) 的集合 E，其中 x 属于诸 \(\mathfrak{g}^{\lambda}(\mathfrak{h})\) 之并）。
\end{enumerate}

\exercisesection{\S~3}

\begin{enumerate}
\def\labelenumi{\arabic{enumi})}
\item
  设 g 是李代数，\(g'\) 是 g 的嘉当子代数。则命题 3 的条件满足。但 \(g'\) 的元素即使在 \(g'\) 中正则，也未必在 g 中正则。
\item
  设 g 是 n 维实李代数，U（相应地 H）是 g 的正则元集合（相应地嘉当子代数集合），\(\text{Int}(\mathfrak{g})\) 是 g 的内自同构群（第三章 §6 第 2 节 定义 2）。
\end{enumerate}

\begin{enumerate}
\def\labelenumi{\alph{enumi})}
\item
  证明：若 \(x\) 与 \(y\) 属于 \(\mathrm{U}\) 的同一连通分支，则 \(\mathfrak{g}^0(x)\) 与 \(\mathfrak{g}^0(y)\) 在 \(\operatorname{Int}(\mathfrak{g})\) 下共轭。
\item
  证明 U 的连通分支个数有限，且该个数被一个只依赖于 n 的常数 \(c(n)\) 所界（应用附录 II 习题 2）。
\item
  推出 \(\operatorname{Int}(\mathfrak{g})\) 在 \(\mathrm{H}\) 上的轨道数 \(\leq c(n)\) 。
\end{enumerate}

\begin{enumerate}
\def\labelenumi{\arabic{enumi})}
\setcounter{enumi}{2}
\tightlist
\item
  设 \(\mathfrak{g}\) 是实李代数，\(\mathfrak{r}\) 是 \(\mathfrak{g}\) 的根基，\(\mathfrak{h}\) 与 \(\mathfrak{h}'\) 是 \(\mathfrak{g}\) 的嘉当子代数，\(\varphi\) 是从 \(\mathfrak{g}\) 到 \(\mathfrak{g} / \mathfrak{r}\) 的典则同态。下列条件等价：
\end{enumerate}

\begin{enumerate}
\def\labelenumi{(\roman{enumi})}
\item
  \(\mathfrak{h}\) 与 \(\mathfrak{h}'\) 在 \(\operatorname{Int}(\mathfrak{g})\) 下共轭；
\item
  \(\varphi(\mathfrak{h})\) 与 \(\varphi(\mathfrak{h}')\) 在 \(\operatorname{Int}(\mathfrak{g}/\mathfrak{r})\) 下共轭。（模仿命题 5 的证明。）
\end{enumerate}

\begin{enumerate}
\def\labelenumi{\arabic{enumi})}
\setcounter{enumi}{3}
\item
  设 \(\mathfrak{g} = \mathfrak{sl}(2, \mathbf{R})\) ，\(x = \begin{pmatrix} 1 & 0 \\ 0 & -1 \end{pmatrix}\) ，\(y = \begin{pmatrix} 0 & -1 \\ 1 & 0 \end{pmatrix}\) 。证明 \(\mathbf{R}x\) 与 \(\mathbf{R}y\) 是 \(\mathfrak{g}\) 的嘉当子代数，且它们在 \(\operatorname{Aut}(\mathfrak{g})\) 下不共轭。
\item
  \begin{enumerate}
  \def\labelenumii{\alph{enumii})}
  \tightlist
  \item
    证明存在李代数 \(\mathfrak{g}\) （ \(k\) 上的），具有基 \((x,y,z,t)\) ，使得
  \end{enumerate}
\end{enumerate}

\[
[ x, y ] = z, [ x, t ] = t, [ y, t ] = 0, [ \mathfrak {g}, z ] = 0.
\]

证明 g 可解，且 \(k = kx + ky + kz\) 是 g 的子代数。

\begin{enumerate}
\def\labelenumi{\alph{enumi})}
\setcounter{enumi}{1}
\item
  证明 g 的初等自同构恰是形如 \(1 + \lambda \operatorname{ad}_{\mathfrak{g}} y + \mu \operatorname{ad}_{\mathfrak{g}} t\) 的映射，其中 \(\lambda, \mu \in k\) 。
\item
  证明 \(1 + ad_{k}x\) 是 k 的初等自同构，它不能开拓为 g 的初等自同构。
\end{enumerate}

\(d\) ) 设 \(\mathfrak{s}\) 是李代数 \(\mathfrak{a}\) 的半单子代数。证明 \(\mathfrak{s}\) 的每个初等自同构都可开拓为 \(\mathfrak{a}\) 的初等自同构。

\begin{enumerate}
\def\labelenumi{\arabic{enumi})}
\setcounter{enumi}{5}
\item
  约化李代数 \(\mathfrak{g}\) 的每个元素都含于一个维数为 \(\mathrm{rk}(\mathfrak{g})\) 的交换子代数中。
\item
  设 \(\mathfrak{g}\) 是李代数，\(\mathfrak{g}'\) 是 \(\mathfrak{g}\) 的子代数且在 \(\mathfrak{g}\) 中约化。设 \(\mathfrak{a}\) 是 \(\mathfrak{g}'\) 的嘉当子代数。证明存在 \(\mathfrak{g}\) 的一个嘉当子代数包含 \(\mathfrak{a}\) （利用 §2 命题 10）。推出 \(\mathrm{rk}(\mathfrak{g}') \leq \mathrm{rk}(\mathfrak{g})\) ，且等号成立当且仅当 \(\mathfrak{g}'\) 具有命题 3 的性质 (i)、(ii)、(iii)。
\item
  设 \(\mathfrak{g}\) 是李代数，\(\mathfrak{a}\) 是 \(\mathfrak{g}\) 的理想，\(\mathfrak{h}\) 是 \(\mathfrak{g}\) 的嘉当子代数，\(\mathcal{C}_{\infty}\mathfrak{g}\) 是 \(\mathfrak{g}\) 的升中心列之并。证明 \(\mathfrak{a} \subset \mathfrak{h}\) 蕴涵 \(\mathfrak{a} \subset \mathcal{C}_{\infty}\mathfrak{g}\) （换言之，\(\mathcal{C}_{\infty}\mathfrak{g}\) 是 \(\mathfrak{g}\) 中含于 \(\mathfrak{h}\) 的最大理想）。（化归到 \(k\) 代数闭的情形，并注意 \(\mathfrak{a}\) 在 \(\mathfrak{g}\) 的每个初等自同构下稳定；于是关系 \(\mathfrak{a} \subset \mathfrak{h}\) 蕴涵 \(\mathfrak{a}\) 含于 \(\mathfrak{g}\) 的每个嘉当子代数；借助 §2 习题 15 得出结论。）
\end{enumerate}

¶ 9) 设 g 是李代数，h 是 g 的嘉当子代数，x 是 h 的元素。设 \(g = h \oplus g^{+}\) 是 g 关于 h 的 Fitting 分解（§1 第 1 节）。

\begin{enumerate}
\def\labelenumi{\alph{enumi})}
\item
  设 \(\mathfrak{n}\) 是最大的半单 \(\mathfrak{h}\) -子模，它含于 \(\mathfrak{g}^0 (x)\cap \mathfrak{g}^+\) 。证明 \(\mathfrak{n} = 0\) 当且仅当 \(\mathfrak{g}^0 (x) = \mathfrak{h}\) ，即 \(x\) 在 \(\mathfrak{g}\) 中正则。
\item
  证明 \(\mathfrak{h} \oplus \mathfrak{n}\) 是 \(\mathfrak{g}\) 的子代数。若 \(\mathfrak{h}'\) 是 \(\mathfrak{h}\) 与 \(\mathfrak{n}\) 的交换子之交，证明 \(\mathfrak{h}'\) 是 \(\mathfrak{h} \oplus \mathfrak{n}\) 的理想，它包含 \(\mathcal{D}\mathfrak{h}\) 与 \(x\) 。由此得出（习题 8）\(\mathfrak{h}' \subset \mathcal{C}_{\infty}(\mathfrak{h} \oplus \mathfrak{n})\) ，从而 \(x \in \mathcal{C}_{\infty}(\mathfrak{h} \oplus \mathfrak{n})\) ，且 \(x\) 属于 \(\mathfrak{h} \oplus \mathfrak{n}\) 的每个嘉当子代数。
\item
  若 \(n \neq 0\) ，则 \(h \oplus n\) 不幂零且有无穷多个嘉当子代数（§2 习题 10）。由此得出 x 属于 g 的无穷多个嘉当子代数。
\end{enumerate}

\(d\) ) \(\mathfrak{g}\) 的元素是正则元，当且仅当它属于唯一的嘉当子代数。

¶ 10) 设 g 是李代数，r 是其根基，n 是其最大幂零理想，a 是其一个 Levi 子代数。

\begin{enumerate}
\def\labelenumi{\alph{enumi})}
\item
  置 \(\mathfrak{g}' = \mathfrak{n} + \mathcal{D}\mathfrak{g}\) 。证明 \(\mathfrak{g}' = \mathfrak{n} \oplus \mathfrak{a}\) 。（利用 \([\mathfrak{g}, \mathfrak{r}]\) 含于 \(\mathfrak{n}\) 这一事实。）若 \(\mathfrak{g} \neq 0\) ，则 \(\mathfrak{g}' \neq 0\) 。
\item
  假设 k 代数闭。设 \((\mathrm{V}_{i})_{i\in\mathrm{I}}\) 是 g-模 g（取伴随表示）的 Jordan-Hölder 序列的商。若 \(x\in r\) ，证明 \(x_{V_{i}}\) 是位似，且 \(x_{V_{i}}=0\) 对一切 i 成立当且仅当 x 属于 n。推出元素 \(y\in g\) 属于 \(g'\) 当且仅当 \(\operatorname{Tr}(y_{\mathrm{V}_{i}})=0\) 对一切 \(i\in I\) 成立。
\item
  用 N 表示 g 中由使 ad x 幂零的元素 x 生成的向量子空间。证明 N 是 g 的子代数（利用 N 在 \(\operatorname{Aut}_{e}(\mathfrak{g})\) 下稳定这一事实）。利用 b) 证明 \(N \subset g'\) 。
\end{enumerate}

\(d\) ) 设 \(\mathfrak{h}\) 是 \(\mathfrak{g}\) 的嘉当子代数。假设存在子集 \(R\) （\(\mathfrak{h}^*\) 中的），使得

\[
\mathfrak {g} = \mathfrak {h} \oplus \bigoplus_ {\alpha \in R} \mathfrak {g} ^ {\alpha} (\mathfrak {h}),
\]

这一假定特别在 k 代数闭时满足。证明 N 此时包含诸 \(\mathfrak{g}^{\alpha}(\mathfrak{h})\) 、Dh 与 n；推出 N 包含 \(g'\) ，故 \(N = g'\) 。

\begin{enumerate}
\def\labelenumi{\alph{enumi})}
\setcounter{enumi}{4}
\tightlist
\item
  若 \(k\) 代数闭且 \(\mathfrak{g} \neq 0\) ，则 \(\mathfrak{g}\) 含有元素 \(x \neq 0\) 使得 \(\operatorname{ad} x\) 幂零。（事实上，此时有 \(\mathfrak{g}' \neq 0\) 。）
\end{enumerate}

¶ 11) 设 g 是李代数，r 是其根基。

\begin{enumerate}
\def\labelenumi{\alph{enumi})}
\item
  设 \(\mathfrak{s}\) 是 \(\mathfrak{g}\) 的 Levi 子代数，\(\mathfrak{k}\) 是 \(\mathfrak{s}\) 的嘉当子代数。证明 \(\mathfrak{k}\) 含于一个嘉当子代数 \(\mathfrak{h}\) （\(\mathfrak{g}\) 中的），后者是 \(\mathfrak{k}\) 与 \(\mathfrak{r}\) 的一个子代数之和。（利用 §2 定理 2、命题 10 及定理 1 的推论 2。）
\item
  设 \(\mathfrak{h}'\) 是 \(\mathfrak{g}\) 的嘉当子代数。证明存在 Levi 子代数 \(\mathfrak{s}'\) （\(\mathfrak{g}\) 的），使得 \(\mathfrak{h}'\) 是 \(\mathfrak{s}'\) 的一个嘉当子代数与 \(\mathfrak{r}\) 的一个子代数之和。（子代数 \(\mathfrak{s}, \mathfrak{k}, \mathfrak{h}\) （\(a\) ) 中的）可以选得使 \(\mathfrak{h} + \mathfrak{r} = \mathfrak{h}' + \mathfrak{r}\) 。置 \(\mathfrak{a} = \mathfrak{h} + \mathfrak{r}\) ，它是可解的。由定理 3，存在 \(x \in \mathcal{C}^{\infty}(\mathfrak{a})\) 使得 \(e^{\mathrm{ad}_{\mathfrak{a}}x}\mathfrak{h} = \mathfrak{h}'\) 。于是 \(e^{\mathrm{ad}_{\mathfrak{g}}x}\) 是 \(\mathfrak{g}\) 的特殊自同构，它把 \(\mathfrak{s}\) 变为所需的 Levi 子代数。）
\item
  设 \(\mathfrak{s}\) 是 \(\mathfrak{g}\) 的 Levi 子代数。设 \(\mathfrak{h}\) 是 \(\mathfrak{g}\) 的嘉当子代数，它是嘉当子代数 \(\mathfrak{k}\) （\(\mathfrak{s}\) 的）与子代数 \(\mathfrak{l}\) （\(\mathfrak{r}\) 的）之和。设 \(\mathfrak{c}\) 是 \(\mathfrak{k}\) 在 \(\mathfrak{r}\) 中的交换子。证明 \(\mathfrak{l}\) 是 \(\mathfrak{c}\) 的嘉当子代数。（对 \(x \in \mathfrak{k}\) ，\(\operatorname{ad}_{\mathfrak{g}} x\) 半单，而 \(\operatorname{ad}_{\mathfrak{h}} x\) 幂零，故 \([\mathfrak{k}, \mathfrak{h}] = 0\) 。若 \(y \in \mathfrak{c}\) 使得 \([y, \mathfrak{l}] \subset \mathfrak{l}\) ，则 \([y, \mathfrak{h}] \subset \mathfrak{h}\) ，故 \(y \in \mathfrak{h} \cap \mathfrak{r} = \mathfrak{l}\) 。）
\item
  设 \(\mathfrak{s}\) 是 \(\mathfrak{g}\) 的 Levi 子代数，\(\mathfrak{k}\) 是 \(\mathfrak{s}\) 的嘉当子代数，\(\mathfrak{c}\) 是 \(\mathfrak{k}\) 在 \(\mathfrak{r}\) 中的交换子，\(\mathfrak{l}\) 是 \(\mathfrak{c}\) 的嘉当子代数。则 \(\mathfrak{h} = \mathfrak{k} + \mathfrak{l}\) 是 \(\mathfrak{g}\) 的嘉当子代数。（设 \(x = y + z\) （ \(y \in \mathfrak{s}, z \in \mathfrak{r}\) ）是正规化子 \(\mathfrak{u}\) （\(\mathfrak{h}\) 在 \(\mathfrak{g}\) 中的）的元素。证明 \([y, \mathfrak{k}] \subset \mathfrak{k}\) ，故 \(y \in \mathfrak{k}\) 且 \(z \in \mathfrak{u}\) 。再证明 \([z, \mathfrak{k}] \subset \mathfrak{h} \cap \mathfrak{r} \subset \mathfrak{c}\) ，故 \([\mathfrak{k}, [\mathfrak{k}, z]] = 0\) ，从而 \([\mathfrak{k}, z] = 0\) 且 \(z \in \mathfrak{c}\) 。最后，\([z, \mathfrak{l}] \subset \mathfrak{l}\) ，故 \(z \in \mathfrak{l}\) 且 \(x \in \mathfrak{h}\) 。）
\item
  设 \(\mathfrak{s},\mathfrak{k},\mathfrak{c}\) 如 \(d\) ) 中所示，\(\mathfrak{q}=[\mathfrak{g},\mathfrak{r}]\) 是 \(\mathfrak{g}\) 的幂零根基。设 \(x\in\mathfrak{q}\) ，\(u\) 是特殊自同构 \(e^{\mathrm{ad}x}\) 。若 \(u(\mathfrak{k})\subset\mathfrak{k}+\mathfrak{c}\) ，则 \(x\in\mathfrak{c}\) 。（考察 \(\rho\) ，即 \(\mathfrak{s}\) 在 \(\mathfrak{q}\) 上的伴随表示，设 \(\mathfrak{q}_{i}\) 是 \(\mathcal{C}^{i+1}\mathfrak{q}\) 在 \(\mathcal{C}^{i}\mathfrak{q}\) 中的在 \(\rho\) 下稳定的补；设 \(\rho_{i}\) 是 \(\rho\) 的由 \(\mathfrak{q}_{i}\) 定义的子表示。设 \(\sigma_{i}=\rho_{i}|\mathfrak{k}\) 。设 \(\mathfrak{q}_{i}^{\prime}\) 是 \(\mathfrak{k}\) 在 \(\mathfrak{q}_{i}\) 中的交换子，\(\mathfrak{q}_{i}^{\prime\prime}\) 是 \(\mathfrak{q}_{i}^{\prime}\) 在 \(\mathfrak{q}_{i}\) 中的在 \(\sigma_{i}\) 下稳定的补。设 \(x=x_{1}^{\prime}+x_{1}^{\prime\prime}+\cdots+x_{n}^{\prime}+x_{n}^{\prime\prime}\) ，其中 \(x_{i}^{\prime}\in\mathfrak{q}_{i}^{\prime},x_{i}^{\prime\prime}\in\mathfrak{q}_{i}^{\prime\prime}\) 。用反证法，设诸 \(x_{i}^{\prime\prime}\) 不全为零，例如 \(x_{1}^{\prime\prime}=\cdots=x_{p-1}^{\prime\prime}=0\) ，\(x_{p}^{\prime\prime}\neq0\) 。若 \(h\in\mathfrak{k}\) ，则 \(u(h)=h+[x_{p}^{\prime\prime},h]+y\) ，其中 \(y\in\mathcal{C}^{p+1}\mathfrak{q}\) 。由于 \(u(h)\in\mathfrak{k}+\mathfrak{c}\) ，这给出 \([x_{p}^{\prime\prime},h]+y\in\mathfrak{q}_{p}^{\prime}+\mathfrak{q}_{p+1}^{\prime}+\cdots+\mathfrak{q}_{n}^{\prime}\) ，从而 \([h,x_{p}^{\prime\prime}]\in\mathfrak{q}_{p}^{\prime}\) ，故 \([h,x_{p}^{\prime\prime}]=0\) 。于是 \(x_{p}^{\prime\prime}=0\) ，矛盾。）
\end{enumerate}

\(f\) ) 设 \(\mathfrak{h}\) 是 \(\mathfrak{g}\) 的嘉当子代数。则 \(\mathfrak{h}\) 可以唯一地表示为 \(\mathfrak{h} \cap \mathfrak{r}\) 与 \(\mathfrak{g}\) 的某个 Levi 子代数的嘉当子代数之和。（关于唯一性，利用 \(e\) ) 及第一章 §6 第 8 节定理 5。）所述 Levi 子代数一般并不唯一。

\(g\) ) 设 \(\mathfrak{h}\) 是 \(\mathfrak{g}\) 的嘉当子代数，并设 \(\mathfrak{t} = \mathfrak{g}^0 (\mathfrak{h} \cap \mathfrak{r})\) 。则 \(\mathfrak{h}\) 是 \(\mathfrak{t}\) 的嘉当子代数。我们有 \(\mathfrak{g} = \mathfrak{t} + \mathfrak{r}\) （对 \(\mathfrak{h} \cap \mathfrak{r}\) 在 \(\mathfrak{g}\) 上的伴随表示用 Fitting 分解）。代数 \(\mathfrak{t} \cap \mathfrak{r}\) 是 \(\mathfrak{t}\) 的根基且幂零（利用 §2 习题 5）。

\begin{enumerate}
\def\labelenumi{\alph{enumi})}
\setcounter{enumi}{7}
\tightlist
\item
  假设 k 代数闭。设 h 是 g 的嘉当子代数。存在 g 的 Levi 子代数 s，使得对一切 \(\lambda \in h^{*}\) ，
\end{enumerate}

\[
\mathfrak {g} ^ {\lambda} (\mathfrak {h}) = \left(\mathfrak {g} ^ {\lambda} (\mathfrak {h}) \cap \mathfrak {s}\right) + \left(\mathfrak {g} ^ {\lambda} (\mathfrak {h}) \cap \mathfrak {r}\right).
\]

（沿用 g) 的记号，把 t 的一个满足 \(\mathfrak{h} = (\mathfrak{h} \cap \mathfrak{s}) + (\mathfrak{h} \cap \mathfrak{r})\) 的 Levi 子代数取作 s；由 b) 它存在。）\(^{7}\)

¶ 12) a) 设 g 是可解李代数，G 是 Aut(g) 的有限子群（若 k = R 或 C，相应地取紧子群）。证明存在 g 的在 G 下稳定的嘉当子代数。（对 dim g 作归纳论证，并化归到如下情形：g 是一个幂零代数 g/n 借助一个交换理想 n 的扩张，其中 n 是非平凡单 g/n-模（参见定理 3 的证明）。此时 g 的嘉当子代数构成一个依附于 n 的仿射空间，G 在其上作用。用重心论证得出结论。）

\begin{enumerate}
\def\labelenumi{\alph{enumi})}
\setcounter{enumi}{1}
\tightlist
\item
  设 g 是可解李代数，s 是 \(\text{Der}(\mathfrak{g})\) 的李子代数。假设 s-模 g 半单。证明存在 g 的在 s 下稳定的嘉当子代数。（同一方法。）
\end{enumerate}

¶ 13) 设 g 是李代数，G 是 Aut(g) 的有限子群。假设 G 超可解（《代数学》第一章 §6 习题 26）。证明存在 g 的在 G 下稳定的嘉当子代数。

（对 dim g 作归纳论证。利用习题 12，化归到 g 半单的情形。若 G ≠ \{1\} ，选取 G 的一个素数阶循环正规子群 C（《代数学》，同前引）。由在 C 下不变的元素组成的子代数 s 在 g 中约化（§1 第 5 节），且异于 g。由归纳假设，s 有一个在 G 下稳定的嘉当子代数 a。我们有 s ≠ 0（第一章 §4 习题 21 c)），故 a ≠ 0。a 在 g 中的交换子 z 异于 g 且在 G 下稳定；选取 z 的一个在 G 下稳定的嘉当子代数 h，并证明 h 是 g 的嘉当子代数，参见命题 3 的推论。）

构造 \(\mathfrak{sl}(2,\mathbf{C})\) 的一个有限自同构群，它同构于 \(A_{4}\) （从而可解），且不保持任何嘉当子代数稳定。

\begin{enumerate}
\def\labelenumi{\arabic{enumi})}
\setcounter{enumi}{13}
\item
  *证明超可解有限群 G 的每个不可约复（相应地实）线性表示由 G 的一个子群的 1 次（相应地 1 次或 2 次）表示所诱导。（对李代数 \(\mathfrak{gl}(n,\mathbf{C})\) 与 \(\mathfrak{gl}(n,\mathbf{R}).\) 应用习题 13。）
\item
  假设 k 代数闭。设 g 是李代数，h 是 g 的嘉当子代数，A 是 g 的子集。假设 A 在 g 中稠密（按 Zariski 拓扑）且在 \(\operatorname{Aut}_{e}(\mathfrak{g})\) 下稳定。证明 \(A \cap h\) 在 h 中稠密。（设 X 是 \(A \cap h\) 的闭包，\(U = h - X\) 。假设 \(U \neq \varnothing\) 。用引理 2 的记号，\(U \times \mathfrak{g}^{\lambda_{1}}(\mathfrak{h}) \times \cdots \times \mathfrak{g}^{\lambda_{p}}(\mathfrak{h})\) 在 F 下的像包含 g 的一个非空开子集。由于该像含于 g-A，这与 A 在 g 中稠密矛盾。）
\item
  设 V 是有限维 k-向量空间，g 是 \(\mathfrak{gl}(V)\) 的李子代数。我们将证明下列三个性质等价：(i) g 的嘉当子代数是交换的，且由半单元素组成。
\end{enumerate}

\begin{enumerate}
\def\labelenumi{(\roman{enumi})}
\setcounter{enumi}{1}
\item
  \(\mathfrak{g}\) 的每个正则元都是半单的。
\item
  g 的半单元素在 g 中按 Zariski 拓扑稠密。
\end{enumerate}

\begin{enumerate}
\def\labelenumi{\alph{enumi})}
\item
  证明 (i) \(\Longrightarrow\) (ii) \(\Longrightarrow\) (iii)。
\item
  设 A 是 \(\mathfrak{g}\) 的半单元素之集。证明 A 在 \(\mathrm{Aut}_e(\mathfrak{g})\) 下稳定。
\item
  证明 (iii) \(\Longrightarrow\) (i)。（（附录 I 习题 1）化归到 k 代数闭的情形。若 h 是 g 的嘉当子代数，利用习题 15 证明 \(A \cap h\) 在 h 中稠密；由于 \([x, y] = 0\) 对 \(x \in A \cap h, y \in h\) 成立，推出 h 交换，由此 (i) 立得。）
\item
  假设 k 是 R、C 或特征零的非离散完备超度量域。给 g 赋予由 k 的拓扑所定义的拓扑。证明性质 (i)、(ii)、(iii) 等价于下列性质：
\end{enumerate}

\begin{enumerate}
\def\labelenumi{(\roman{enumi})}
\setcounter{enumi}{3}
\tightlist
\item
  \(\mathfrak{g}\) 的半单元素在 \(\mathfrak{g}\) 中稠密。
\end{enumerate}

（证明 (iv) \(\Longrightarrow\) (iii) 与 (ii) \(\Longrightarrow\) (iv)，参见附录 I 习题 4。）

\begin{enumerate}
\def\labelenumi{\arabic{enumi})}
\setcounter{enumi}{16}
\tightlist
\item
  假设 \(k\) 代数闭。设 \(\mathfrak{g}\) 是李代数，\(\mathfrak{h}\) 是其嘉当子代数，A 是 \(\mathfrak{h}\) 的中心的子集。用 \(\mathrm{E}_{\mathfrak{g}}\) 表示 \(\operatorname{Aut}(\mathfrak{g})\) 的第 2 节中记作 E 的子群。证明：若 \(s\) 是 \(\operatorname{Aut}(\mathfrak{g})\) 的元素且 \(s\mathrm{A} = \mathrm{A}\) ，则存在 \(t \in \mathrm{E}_{\mathfrak{g}}\) 使得 \(t\mathfrak{h} = \mathfrak{h}\) 且 \(t|\mathrm{A} = \mathrm{Id}_{\mathrm{A}}\) ；特别地，\(ts|_{\mathrm{A}} = s|_{\mathrm{A}}\) 。（设 \(\mathfrak{a}\) 是 A 在 \(\mathfrak{g}\) 中的交换子；由于 \(\mathfrak{h}\) 与 \(sh\) 都是 \(\mathfrak{a}\) 的嘉当子代数，存在 \(\theta \in \mathrm{E}_{\mathfrak{a}}\) 使得 \(\theta(s\mathfrak{h}) = \mathfrak{h}\) ；把 \(t\) 取为 \(\mathrm{E}_{\mathfrak{g}}\) 中开拓 \(\theta\) 的一个元素。）
\end{enumerate}

¶ 18) 设 g 是李代数，h 是 g 的嘉当子代数，Ug（相应地 Uh）是 g（相应地 h）的包络代数。Ug 上的线性形式 \(\varphi\) 称为中心的，若它在 {[}Ug, Ug{]} 上为零，即若 \(\varphi(a.b) = \varphi(b.a)\) 对一切 \(a, b \in Ug\) 成立。

\(a)\) 设 \(x, y \in \mathfrak{g}\) 。假设存在 \(s \in \operatorname{Aut}_e(\mathfrak{g})\) 使得 \(s(x) = y\) 。证明

\[
\varphi (x ^ {n}) = \varphi (y ^ {n})
\]

对一切 \(n \in \mathbf{N}\) 与每个中心线性形式 \(\varphi\) （\(\mathrm{Ug}\) 上的）成立。

\begin{enumerate}
\def\labelenumi{\alph{enumi})}
\setcounter{enumi}{1}
\item
  设 \(\varphi\) 是 Ug 上的中心线性形式，它在 Uh 上的限制为零。证明 \(\varphi = 0\) 。（可以假设 k 代数闭。由 a) 推出此时 \(\varphi(x^{n}) = 0\) 对一切 \(n \in N\) 及一切正则的 \(x \in g\) 成立；用稠密性论证去掉正则性假设。）
\item
  证明 \(Ug = [Ug, Ug] + Uh\) 。
\end{enumerate}

\(d\) ) 设 V 是半单 \(\mathfrak{g}\) -模。证明 V 作为 \(\mathfrak{h}\) -模半单。特别地，\(V^{\lambda}(\mathfrak{h}) = V_{\lambda}(\mathfrak{h})\) 对一切 \(\lambda \in \mathfrak{h}^{*}\) 成立。

\begin{enumerate}
\def\labelenumi{\alph{enumi})}
\setcounter{enumi}{4}
\tightlist
\item
  设 \(V'\) 是半单 g-模。假设 V 与 \(V'\) 作为 h-模同构。证明它们作为 g-模同构。（若 \(a \in Uh\) ，注意 \(\operatorname{Tr}(a_{\mathrm{V}}) = \operatorname{Tr}(a_{\mathrm{V}'})\) 。）利用 b) 或 c) 推出 \(\operatorname{Tr}(x_{\mathrm{V}}) = \operatorname{Tr}(x_{\mathrm{V}'})\) 对一切 \(x \in Ug\) 成立，并借助《代数学》第八章得出结论。
\end{enumerate}

若 k 代数闭，则假定“V 与 V' 是 h-同构的”等价于说 \(\dim\mathrm{V}_{\lambda}(\mathfrak{h})=\dim\mathrm{V}_{\lambda}^{\prime}(\mathfrak{h})\) 对一切 \(\lambda\in\mathfrak{h}^{*}\) 成立。

\exercisesection{\S~4}

记号与假定同 §4 第 1、2、3 节。

\begin{enumerate}
\def\labelenumi{\arabic{enumi})}
\tightlist
\item
  取 \(\mathrm{G} = \mathbf{GL}_{n}(k), n \geq 0.\)
\end{enumerate}

\begin{enumerate}
\def\labelenumi{\alph{enumi})}
\item
  证明 \(r_{\mathrm{Ad}}^0 (g) = n\) 对一切 \(g \in \mathrm{G}\) 成立。
\item
  证明元素 \(g \in G\) 正则，当且仅当其特征多项式 \(\mathrm{P}_{g}(\mathrm{T}) = \det(\mathrm{T} - g)\) 可分；这等价于说 \(\mathrm{P}_{g}(\mathrm{T})\) 的判别式（《代数学》第四章 §1 第 10 节） \(\neq 0\) 。
\end{enumerate}

\begin{enumerate}
\def\labelenumi{\arabic{enumi})}
\setcounter{enumi}{1}
\item
  构造一个李群 G，使函数 \(r_{Ad}^{0}\) 不是常数。（取 g 交换且 \(\neq 0\) ，且 Ad 非平凡。）
\item
  设 \((\rho_{i})_{i\in I}\) 是 G 的可数族解析线性表示。证明 G 中对所有 \(\rho_{i}\) 都正则的元素构成 G 的稠密子集。构造一个例子说明可数性假定不可省去。
\item
  假设 \(k = \mathbf{C}\) 且 \(\mathbf{G}\) 连通。证明下列性质等价：
\end{enumerate}

\begin{enumerate}
\def\labelenumi{(\roman{enumi})}
\item
  G 幂零。
\item
  每个 \(\neq 1\) 的 \(G\) 中元素都正则。
\end{enumerate}

（先证明 (ii) 蕴涵

\begin{enumerate}
\def\labelenumi{(\roman{enumi})}
\setcounter{enumi}{1}
\tightlist
\item
  \(^{\prime}\) g 的每个 \(\neq 0\) 的元素都正则。
\end{enumerate}

再注意，若 \(g \neq 0\) ，则 g 含有元素 \(x \neq 0\) 使 ad x 幂零，参见 §3 习题 10。推出 g 幂零，从而得 (i)。

\exercisesection{\S~5}

\begin{enumerate}
\def\labelenumi{\arabic{enumi})}
\item
  证明第一章 §5 习题 6 中所考虑的可解李代数不同构于任何可分解李代数。
\item
  设 u（相应地 v）是 V 的非零半单（相应地幂零）自同态。则映射 \(\lambda u \mapsto \lambda v (\lambda \in k)\) 是从 g = ku 到 \(g' = kv\) 的同构，但它不把半单元素映为半单元素。
\item
  设 \(u\) 是 \(V\) 的既非半单亦非幂零的自同态。则 \(\mathfrak{g} = ku\) 不可分解，但 \(\operatorname{ad}_{\mathfrak{g}}\mathfrak{g}\) 可分解。
\end{enumerate}

¶ 4) 设 g 是 gl(V) 的可分解李子代数。设 q 是 g 的李子代数，其恒同表示半单。存在 \(a \in \operatorname{Aut}_{e}(\mathfrak{g})\) 使 \(a(\mathfrak{q})\) 含于命题 7 的子代数 m。（模仿命题 7 (ii) 的证明。）

¶ 5) 设 g 是 gl(V) 的李子代数。若对一切 \(x \in \mathfrak{g}\) ，\(x\) 的复本（第一章 §5 习题 14）都属于 g，则称 g 为代数的。这样的代数是可分解的。

\begin{enumerate}
\def\labelenumi{\alph{enumi})}
\tightlist
\item
  用 \(a(\mathfrak{g})\) 表示 \(\mathfrak{gl}(\mathrm{V})\) 中包含 \(\mathfrak{g}\) 的最小代数子代数。则
\end{enumerate}

\[
a (\mathfrak {g}) \supset e (\mathfrak {g}) \supset \mathfrak {g}.
\]

给出一个 \(a(\mathfrak{g})\) 与 \(e(\mathfrak{g})\) 相异的例子（取 V 为 2 维，\(\mathfrak{g}\) 为 1 维）。

\begin{enumerate}
\def\labelenumi{\alph{enumi})}
\setcounter{enumi}{1}
\item
  证明：若 \(\mathfrak{n}\) 是 \(\mathfrak{g}\) 的理想，则 \(\mathfrak{n}\) 与 \(a(\mathfrak{n})\) 都是 \(a(\mathfrak{g})\) 的理想，且 \([a(\mathfrak{g}), a(\mathfrak{n})] = [\mathfrak{g}, \mathfrak{n}]\) （模仿命题 4 的证明）。推出 \(\mathcal{D}^i a(\mathfrak{g}) = \mathcal{D}^i \mathfrak{g}\) 对 \(i \geq 1\) 成立，且 \(\mathcal{C}^i a(\mathfrak{g}) = \mathcal{C}^i \mathfrak{g}\) 对 \(i \geq 2\) 成立。
\item
  证明由幂零元素组成的李代数都是代数的。\(^{8}\)
\end{enumerate}

¶ 6) 设 g 是 \(\mathfrak{gl}(V)\) 的半单子代数，\(\mathbf{T}(V) = \bigoplus_{n=0}^{\infty} \mathbf{T}^{n}(V)\) 是 V 的张量代数，\(\mathbf{T}(V)^{\mathfrak{g}}\) 是 \(\mathbf{T}(V)\) 中在 g 下不变的元素之集（参见第三章附录），\(\overline{g}\) 表示这样的 \(u \in \mathfrak{gl}(V)\) 组成的集合：u.x = 0 对一切 \(x \in \mathbf{T}(V)^{\mathfrak{g}}\) 成立。我们将证明 \(\overline{g} = g\) 。

\begin{enumerate}
\def\labelenumi{\alph{enumi})}
\item
  证明 g 在 V 的对偶 V* 上的表示同构于它在 \(\bigwedge^{p-1}V\) 上的表示，其中 \(p = \dim V\) （利用 g 含于 \(\mathfrak{sl}(V)\) 这一事实）。推出 \(\mathbf{T}_{n,m} = \mathbf{T}^{n}(\mathrm{V}) \otimes \mathbf{T}^{m}(\mathrm{V}^{*})\) 中在 g 下不变的每个元素都在 \(\overline{g}\) 下不变，且 \(\overline{g}\) 是代数的（习题 5）。
\item
  设 W 是某个 \(T_{n,m}\) 的向量子空间。假设 W 在 g 下稳定。证明 W 在 \(\overline{g}\) 下稳定（若 \(e_{1},\ldots,e_{r}\) 是 W 的基，注意 \(e_{1}\wedge\cdots\wedge e_{r}\) 在 g 下不变，从而在 \(\overline{g}\) 下不变）。
\end{enumerate}

推出 g 是 \(\overline{g}\) 的理想，且 \(\overline{g}/g\) 交换（参见命题 4 的证明）。我们有 \(\overline{g} = g \times c\) ，其中 c 是 \(\overline{g}\) 的中心。

\begin{enumerate}
\def\labelenumi{\alph{enumi})}
\setcounter{enumi}{2}
\item
  设 R 是 \(\mathfrak{gl}(V)\) 中由 1 与 g 生成的结合子代数。证明 c 含于 R 的中心（注意 \(\overline{g}\) 含于 R 的双交换子，而后者等于 R）。推出 c 的元素都是半单的。
\item
  设 \(x \in c\) 。证明 x 的复本属于 c（第一章 §5 习题 14）。证明 \(\operatorname{Tr}(sx) = 0\) 对一切 \(s \in g\) 成立；推出 \(\operatorname{Tr}(sx) = 0\) 对一切 \(s \in \overline{g}\) 成立，从而 x 幂零（同前引）。
\item
  证明 \(\mathfrak{c} = 0\) 且 \(\overline{\mathfrak{g}} = \mathfrak{g}\) ，把 \(c\) 与 \(d\) ) 结合起来。
\end{enumerate}

\begin{enumerate}
\def\labelenumi{\arabic{enumi})}
\setcounter{enumi}{6}
\tightlist
\item
  设 \(\mathfrak{g}\) 是 \(\mathfrak{gl}(\mathrm{V})\) 的李子代数。设 \(m, n\) 是两个整数 \(\geq 0\) ，W 与 \(\mathrm{W}'\) 是 \(\mathbf{T}^m(\mathrm{V}) \otimes \mathbf{T}^n(\mathrm{V}^*)\) 的两个向量子空间，其中 \(\mathrm{V}^*\) 是 V 的对偶。假设 \(\mathrm{W}' \subset \mathrm{W}\) ，且 W 与 \(\mathrm{W}'\) 在 \(\mathfrak{g}\) 于 \(\mathbf{T}^m(\mathrm{V}) \otimes \mathbf{T}^n(\mathrm{V}^*)\) 上的自然表示下稳定。证明 W 与 \(\mathrm{W}'\) 此时在 \(e(\mathfrak{g})\) 下稳定。若用 \(\pi\) 表示 \(e(\mathfrak{g})\) 由此在 \(\mathrm{W/W}'\) 上给出的表示，证明 \(\pi e(\mathfrak{g})\) 是 \(\pi(\mathfrak{g})\) 的可分解包络（利用定理 1）。
\end{enumerate}

推出 \(\operatorname{ad} e(\mathfrak{g})\) 是 \(\operatorname{ad} \mathfrak{g}\) 在 \(\mathfrak{gl}(\mathfrak{g})\) 中的可分解包络。

\begin{enumerate}
\def\labelenumi{\arabic{enumi})}
\setcounter{enumi}{7}
\tightlist
\item
  设 g 是 \(\mathfrak{gl}(V)\) 的李子代数，h 是 g 的嘉当子代数。
\end{enumerate}

\begin{enumerate}
\def\labelenumi{\alph{enumi})}
\tightlist
\item
  证明 \(e(\mathfrak{g}) = e(\mathfrak{h}) + \mathcal{D}\mathfrak{g} = e(\mathfrak{h}) + \mathfrak{g}\) 。
\end{enumerate}

（注意 \(e(\mathfrak{h}) + \mathcal{D}\mathfrak{g}\) 可分解（定理 1 的推论 1），包含 \(\mathfrak{g} = \mathfrak{h} + \mathcal{D}\mathfrak{g}\) ，且含于 \(e(\mathfrak{g})\) ；故它就是 \(e(\mathfrak{g})\) 。）

\begin{enumerate}
\def\labelenumi{\alph{enumi})}
\setcounter{enumi}{1}
\item
  我们有 \(e(\mathfrak{h}) \cap \mathfrak{g} = \mathfrak{h}\) （注意 \(e(\mathfrak{h}) \cap \mathfrak{g}\) 幂零）。
\item
  设 \(x\) 是 \(e(\mathfrak{h})\) 在 \(e(\mathfrak{g})\) 中的正规化子的元素。证明 \(x \in e(\mathfrak{h})\) 。（写 \(x = y + z\) ，其中 \(y \in e(\mathfrak{h}), z \in \mathfrak{g}\) ，参见 a)；注意 \([z, \mathfrak{h}] \subset e(\mathfrak{h}) \cap \mathfrak{g} = \mathfrak{h}\) ，故 \(z \in \mathfrak{h}\) 。）
\end{enumerate}

\(d\) ) 证明 \(e(\mathfrak{h})\) 是 \(e(\mathfrak{g})\) 的嘉当子代数。

\begin{enumerate}
\def\labelenumi{\arabic{enumi})}
\setcounter{enumi}{8}
\tightlist
\item
  设 \(\mathfrak{g}\) 是 \(\mathfrak{gl}(\mathrm{V})\) 的李子代数。证明 §3 习题 16 的条件 (i)、(ii)、(iii) 等价于：
\end{enumerate}

\begin{enumerate}
\def\labelenumi{(\alph{enumi})}
\setcounter{enumi}{21}
\tightlist
\item
  \(\mathfrak{g}\) 可分解，且与 \(\mathfrak{g} / \mathfrak{n}_{\mathrm{V}}(\mathfrak{g})\) 有相同的秩。
\end{enumerate}

（若 \(\mathfrak{h}\) 是 \(\mathfrak{g}\) 的嘉当子代数，则条件“\(\mathfrak{g}\) 与 \(\mathfrak{g}/\mathfrak{n}_{\mathrm{V}}(\mathfrak{g})\) 有相同的秩”等价于说 \(\mathfrak{h} \cap \mathfrak{n}_{\mathrm{V}}(\mathfrak{g}) = 0\) ，即 \(\mathfrak{h}\) 没有 \(\neq 0\) 的幂零元素（参见 §2 习题 14）。推出 (i) 与 (v) 等价。）

\begin{enumerate}
\def\labelenumi{\arabic{enumi})}
\setcounter{enumi}{9}
\tightlist
\item
  设 \(k'\) 是 \(k\) 的一个扩张，\(\mathfrak{g}'\) 是下式的 \(k'\) -李子代数
\end{enumerate}

\[
\mathfrak {g l} (\mathrm{V} \otimes_ {k} k ^ {\prime}) = \mathfrak {g l} (\mathrm{V}) \otimes_ {k} k ^ {\prime}.
\]

\begin{enumerate}
\def\labelenumi{\alph{enumi})}
\item
  证明存在 \(\mathfrak{gl}(V)\) 的最小李子代数 g，使得 \(g\otimes_{k}k'\) 包含 \(g'\) 。
\item
  假设 \(k'\) 代数闭，并用 G 表示 \(k'\) 的 k-自同构群；此群以自然方式作用于 \(V \otimes_{k} k'\) 。证明 \(g \otimes_{k} k'\) 是由 \(g'\) 在 G 下的共轭元生成的李子代数（利用 k 是 \(k'\) 中 G-不变元所成的域这一事实）。
\item
  证明 \(\mathfrak{g}\) 可分解，若 \(\mathfrak{g}'\) 可分解。（化归到 \(k'\) 代数闭的情形，并利用 \(b\) ) 以及定理 1 的推论 1 与 3。）
\end{enumerate}

¶ 11) 作为例外，本习题中我们假设 k 是特征为 p \textgreater{} 0 的完满域。

设 g 是李 p-代数（第一章 §1 习题 20）。若 \(x \in g\) ，用 \(\langle x \rangle\) 表示包含 x 的最小李 p-子代数。它是交换的，且作为 k-向量空间由诸 \(x^{p^{i}}\) 生成，其中 \(i = 0, 1, \ldots\) 。称 x 为幂零的（相应地半单的），若 \(\langle x \rangle\) 的 p-映射幂零（相应地双射）。

\begin{enumerate}
\def\labelenumi{\alph{enumi})}
\item
  证明 x 可以唯一地分解为形式 \(x = s + n\) ，其中 \(s, n \in \langle x \rangle\) ，s 半单，n 幂零（应用第一章 §1 习题 23）。若 f 是从 g 到 \(\mathfrak{gl}(\mathrm{V})\) 的 p-同态，则 \(f(s)\) 与 \(f(n)\) 是自同态 \(f(x)\) 的半单分量与幂零分量；这特别适用于 f = ad。
\item
  g 的子代数称为可分解的，若它包含其元素的半单分量与幂零分量。证明：若 b 与 c 是 g 的向量子空间且 \(b \subset c\) ，则 \(x \in g\) 中使得 \([x, c] \subset b\) 的那些 x 组成的集合是可分解的（与命题 3 的证明相同）；特别地，g 的每个嘉当子代数都可分解。
\item
  设 t 是 g 的由半单元素组成的交换子代数，且对此性质极大。设 h 是 t 在 g 中的交换子。设 \(x \in h\) ，\(x = s + n\) 是其典则分解；由于 h 可分解（参见 b)），\(s, n \in h\) 。证明由 t 与 s 生成的子代数交换且由半单元素组成，故与 t 一致。推出 \(ad_{h}x = ad_{h}n\) 幂零，从而 h 幂零。由于 \(\mathfrak{h} = \mathfrak{g}^{0}(\mathfrak{h})\) ，h 是 g 的嘉当子代数（§2 命题 4）。
\end{enumerate}

特别地，有限域上的每个李 p-代数都有嘉当子代数。\(^{9}\)

\exercisesection{附录~I}

用 V 表示 \(k\) 上的有限维向量空间。

\begin{enumerate}
\def\labelenumi{\arabic{enumi})}
\item
  设 \(k'\) 是 k 的扩张，并设 \(\mathrm{V}_{(k')} = \mathrm{V} \otimes_k k'\) 。证明 \(\mathrm{V}_{(k')}\) 上的 Zariski 拓扑在 V 上诱导出 V 的 Zariski 拓扑，且 V 在 \(\mathrm{V}_{(k')}\) 中稠密。
\item
  假设 \(\mathrm{V}\) 是两个向量空间 \(\mathrm{V}_1\) 与 \(\mathrm{V}_2\) 的乘积。
\end{enumerate}

\begin{enumerate}
\def\labelenumi{\alph{enumi})}
\item
  V 上的 Zariski 拓扑细于 \(V_{1}\) 与 \(V_{2}\) 上 Zariski 拓扑的乘积拓扑；若 \(V_{1} \neq 0\) 且 \(V_{2} \neq 0\) ，则严格细于它。
\item
  若 \(A_{1}\) （相应地 \(A_{2}\) ）是 \(V_{1}\) （相应地 \(V_{2}\) ）的子集，则 \(A_{1} \times A_{2}\) 的闭包是 \(A_{1}\) 与 \(A_{2}\) 的闭包之积。
\end{enumerate}

\begin{enumerate}
\def\labelenumi{\arabic{enumi})}
\setcounter{enumi}{2}
\tightlist
\item
  假设 \(k\) 代数闭。设 \(A\) 与 \(B\) 是 \(V\) 的两个闭子集，\(\mathfrak{a}\) （相应地 \(\mathfrak{b}\) ）是由 \(f \in A_V\) 中在 \(A\) （相应地 \(B\) ）上为零者组成的集合。证明下列性质等价：
\end{enumerate}

\begin{enumerate}
\def\labelenumi{(\roman{enumi})}
\item
  A ∩ B = ∅。
\item
  \(\mathfrak{a} + \mathfrak{b} = \mathrm{A_V}\)。
\item
  存在多项式函数 \(f\) （\(V\) 上的），它在 \(A\) 上等于 1，在 \(B\) 上等于 0。
\end{enumerate}

（利用 Hilbert 零点定理（《交换代数》第五章 §3 第 3 节）证明 (i) \(\Longrightarrow\) (ii)。）

\begin{enumerate}
\def\labelenumi{\arabic{enumi})}
\setcounter{enumi}{3}
\tightlist
\item
  假设 \(k\) 是非离散完备赋值域。用 \(\mathcal{T}\) （相应地 \(\mathcal{Z}\) ）表示 V 上的 Banach 空间拓扑（相应地 Zariski 拓扑）。
\end{enumerate}

\begin{enumerate}
\def\labelenumi{\alph{enumi})}
\item
  证明 \(\mathcal{T}\) 细于 \(\mathcal{L}\) （且当 \(V \neq 0\) 时严格细）。
\item
  证明 V 的每个非空 \(\mathcal{Z}\) -开子集都在 \(\mathcal{T}\) 意义下稠密。
\end{enumerate}

\exercisesection{附录~II}

¶ 1) 设 X 是局部连通拓扑空间，\(\mathcal{C}(\mathrm{X})\) 是 X 上实值连续函数组成的空间，\(d\) 是整数 \(\geq 0\) 。设 \(\mathrm{F} \in \mathcal{C}(\mathrm{X})[\mathrm{T}]\) 是次数为 \(d\) 、系数属于 \(\mathcal{C}(\mathrm{X})\) 的首一多项式：

\[
\mathrm{F} = \mathrm{T} ^ {d} + \mathrm{T} ^ {d - 1} f _ {1} + \dots + f _ {d}, \quad f _ {i} \in \mathcal {C} (\mathrm{X}).
\]

通过下式把 F 等同于 \(R \times X\) 上的一个函数：

\[
\mathrm{F} (t, x) = t ^ {d} + t ^ {d - 1} f _ {1} (x) + \dots + f _ {d} (x) \text {if} t \in \mathbf {R}, x \in \mathrm{X}.
\]

设 \(\Delta \in \mathcal{C}(\mathrm{X})\) 是多项式 F 的判别式（《代数学》第四章 §1 第 10 节）。

若 U 是 X 的开子集，用 \(Z_{U}\) 表示由 \((t,x)\) （其中 \(t \in R\) 且 \(x \in U\) ）组成且满足 \(F(t,x) = 0\) 的集合；这是 \(R \times U\) 的闭子集。

\begin{enumerate}
\def\labelenumi{\alph{enumi})}
\item
  证明投影 \(\mathrm{pr}_2: \mathrm{Z}_{\mathrm{U}} \to \mathrm{U}\) 是真的（《一般拓扑学》第一章 §10）。
\item
  假设 U 连通，且 \(\Delta(x) \neq 0\) 对一切 \(x \in U\) 成立。证明 \(Z_{U} \to U\) 是 U 的次数 \(\leq d\) 的覆盖（《一般拓扑学》第十一章），且 \(R \times U - Z_{U}\) 的连通分支个数 \(\leq d + 1\) 。
\item
  设 \(X'\) 是 X 中使 \(\Delta\) \(\neq 0\) 的点组成的集合。假设 \(X'\) 在 X 中稠密。用 \(\mathcal{A}(\text{resp. } \mathcal{B})\) 表示 \(X'\) （相应地 \(R \times X - Z_{X}\) ）的连通分支之集。证明
\end{enumerate}

\[
\operatorname{Card} (\mathscr {B}) \leq (d + 1) \operatorname{Card} (\mathscr {A}) \quad (\text { use   } b)).
\]

\begin{enumerate}
\def\labelenumi{\alph{enumi})}
\setcounter{enumi}{3}
\tightlist
\item
  假设 X 连通，且 \(d \geq 1\) 。证明
\end{enumerate}

\[
\operatorname{Card} (\mathcal {B}) \leq 1 + d \operatorname{Card} (\mathcal {A}).
\]

¶ 2) 设 V 是有限 n 维实向量空间，F 是 V 上的 d 次多项式函数。设 \(V'\) 是 V 中使 \(F \neq 0\) 的点组成的集合。证明 \(V'\) 的连通分支个数有限，且被一个只依赖于 n 与 d 的常数所界。（对 n 作归纳论证。化归到 F 无重因子的情形，并证明 V 可分解为 \(R \times X\) ，使习题 1 的结果可应用于 F。）\(^{10}\)

\chapter{分裂半单李代数}

本章中，k 表示一个特征为 0 的（交换）域。除另有说明外，“向量空间”均指“k 上的向量空间”；“李代数”等同理。

\section{\texorpdfstring{李代数 \(\mathfrak{sl}(2,k)\) 及其表示}{李代数 \textbackslash mathfrak\{sl\}(2,k) 及其表示}}

\subsection{\texorpdfstring{\(\mathfrak{sl}(2,k)\) 的典则基}{\textbackslash mathfrak\{sl\}(2,k) 的典则基}}

引理 1. 设 \(\mathbf{A}\) 是 \(k\) 上的结合代数，\(H\) 与 \(X\) 是 \(\mathbf{A}\) 中满足 \([H, X] = 2X\) 的元素。

\begin{enumerate}
\def\labelenumi{(\roman{enumi})}
\item
  \([H,X^n ] = 2nX^n\) 对任何整数 \(n\geq 0\) 成立
\item
  若 \(Z\) 是 \(A\) 中满足 \([Z, X] = H\) 的元素，则对任何整数 \(n > 0\) ，
\end{enumerate}

\[
[ Z, X ^ {n} ] = n X ^ {n - 1} (H + n - 1) = n (H - n + 1) X ^ {n - 1}.
\]

映射 \(T \mapsto [H, T]\) （从 \(A\) 到 \(A\) 的）是一个导子，由此得 (i)。在 (ii) 的假定下，

\[
\begin{array}{l} [ Z, X ^ {n} ] = \sum_ {i + j = n - 1} X ^ {i} H X ^ {j} \\ \qquad = \sum_ {i + j = n - 1} (X ^ {i} X ^ {j} H + X ^ {i} 2 j X ^ {j}) \\ \qquad = n X ^ {n - 1} H + 2 X ^ {n - 1} \frac {n (n - 1)}{2} \\ \qquad = n X ^ {n - 1} (H + n - 1). \end{array}
\]

另一方面，由 (i) 有 \(X^{n-1}(H+n-1)=(H-n+1)X^{n-1}\) 。证毕。

回顾我们用 \(\mathfrak{sl}(2,k)\) 表示由 2 阶、迹为零、元素在 k 中的方阵组成的李代数。此李代数是 3 维单李代数（第一章 §6 第 7 节 例）。\(\mathfrak{sl}(2,k)\) 的典则基是基 \((X_{+},X_{-},H)\) ，其中

\[
X _ {+} = \left( \begin{array}{c c} 0 & 1 \\ 0 & 0 \end{array} \right) \quad X _ {-} = \left( \begin{array}{c c} 0 & 0 \\ - 1 & 0 \end{array} \right) \quad H = \left( \begin{array}{c c} 1 & 0 \\ 0 & - 1 \end{array} \right).
\]

我们有

\[
[ H, X _ {+} ] = 2 X _ {+} \quad [ H, X _ {-} ] = - 2 X _ {-} \quad [ X _ {+}, X _ {-} ] = - H.\tag{1}
\]

由于 \(\mathfrak{sl}(2,k)\) 的恒同表示是单射，H 是 \(\mathfrak{sl}(2,k)\) 的半单元素，而 \(X_{+}, X_{-}\) 是 \(\mathfrak{sl}(2,k)\) 的幂零元素（第一章 §6 第 3 节 定理 3）。由第七章 §2 第 1 节 例 4，kH 是 \(\mathfrak{sl}(2,k)\) 的嘉当子代数。映射 \(U \mapsto -^{t}U\) 是李代数 \(\mathfrak{sl}(2,k)\) 的对合自同构，称为 \(\mathfrak{sl}(2,k)\) 的典则对合；它把 \((X_{+}, X_{-}, H)\) 变为 \((X_{-}, X_{+}, -H)\) 。

引理 2. 在 \(\mathfrak{sl}(2,k)\) 的包络代数中，

\[
[ H, X _ {+} ^ {n} ] = 2 n X _ {+} ^ {n} \quad [ H, X _ {-} ^ {n} ] = - 2 n X _ {-} ^ {n}
\]

对任何整数 \(n \geq 0\) 成立，且

\[
[ X _ {-}, X _ {+} ^ {n} ] = n X _ {+} ^ {n - 1} (H + n - 1) = n (H - n + 1) X _ {+} ^ {n - 1}
\]

\[
[ X _ {+}, X _ {-} ^ {n} ] = n X _ {-} ^ {n - 1} (- H + n - 1) = n (- H - n + 1) X _ {-} ^ {n - 1}
\]

若 n \textgreater{} 0。

第一与第三个关系式由引理 1 得出。其余的可用 \(\mathfrak{sl}(2,k)\) 的典则对合从它们推出。

\subsection{\texorpdfstring{\(\mathfrak{sl}(2,k)\) -模的本原元}{\textbackslash mathfrak\{sl\}(2,k) -模的本原元}}

设 E 是 \(\mathfrak{sl}(2,k)\) -模。若 \(A \in \mathfrak{sl}(2,k)\) 且 \(x \in E\) ，我们常把 \(A_{E}x\) 写作 Ax。设 \(\lambda \in k\) 。若 Hx = \(\lambda x\) ，则按语言滥用，我们称 x 是 E 中权为 \(\lambda\) 的元素，或称 \(\lambda\) 是 x 的权。若 E 有限维，则 \(H_{E}\) 半单，故权为 \(\lambda\) 的元素之集是 E 相对于 \(H_{E}\) 与 \(\lambda\) 的准素子空间（参见第七章 §1 第 1 节）。

引理 3. 若 \(x\) 是权为 \(\lambda\) 的元素，则 \(X_{+}x\) 是权为 \(\lambda + 2\) 的元素，而 \(X_{-}x\) 是权为 \(\lambda - 2\) 的元素。

事实上，\(HX_{+}x = [H,X_{+}]x + X_{+}Hx = 2X_{+}x + X_{+}\lambda x = (\lambda +2)X_{+}x\) ，同理 \(HX_{-}x = (\lambda -2)X_{-}x\) （另见第七章 §1 第 3 节 命题 10 (ii)）。

定义 1. 设 E 是 \(\mathfrak{sl}(2,k)\) -模。E 的元素称为本原元，如果它是 \(H_{E}\) 的非零特征向量且属于 \(X_{+E}\) 的核。

E 的非零元素 e 是本原元，当且仅当 ke 在 \(kH + kX_{+}\) 的作用下稳定；这例如可由引理 3 得出。

例 元素 \(X_{+}\) 对 \(\mathfrak{sl}(2,k)\) 的伴随表示是权 2 的本原元。元素 \((1,0)\) （\(k^{2}\) 中的）对 \(\mathfrak{sl}(2,k)\) 在 \(k^{2}\) 上的恒同表示是权 1 的本原元。

引理 4. 设 \(\mathbf{E}\) 是非零有限维 \(\mathfrak{sl}(2,k)\) -模。则 \(\mathbf{E}\) 有本原元。

由于 \(X_{+}\) 是 \(\mathfrak{sl}(2,k)\) 的幂零元素，\(X_{+E}\) 幂零。假设 \(X_{+E}^{m-1} \neq 0\) 且 \(X_{+E}^{m} = 0\) 。由引理 2，

\[
m (H _ {\mathrm{E}} - m + 1) X _ {+ \mathrm{E}} ^ {m - 1} = [ X _ {- \mathrm{E}}, X _ {+ \mathrm{E}} ^ {m} ] = 0,
\]

因而 \(X_{+}^{m - 1}(\mathrm{E}) - \{0\}\) 中的元素都是本原的。

命题 1. 设 E 是 \(\mathfrak{sl}(2,k)\) -模，\(e\) 是 E 的权为 \(\lambda\) 的本原元。置 \(e_n = \frac{(-1)^n}{n} X_-^n e\) （对 \(n \geq 0\) ），并置 \(e_{-1} = 0\) 。则

\[
\left\{ \begin{array}{l l} H e _ {n} & = (\lambda - 2 n) e _ {n} \\ X _ {-} e _ {n} & = - (n + 1) e _ {n + 1} \\ X _ {+} e _ {n} & = (\lambda - n + 1) e _ {n - 1}. \end{array} \right.\tag{2}
\]

第一个公式由引理 3 得出，第二个由 \(e_n\) 的定义得出。我们对 \(n\) 作归纳来证明第三个。它在 \(n = 0\) 时成立，因为 \(e_{-1} = 0\) 。若 \(n > 0\) ，

\[
\begin{array}{r l} & n X _ {+} e _ {n} = - X _ {+} X _ {-} e _ {n - 1} = - [ X _ {+}, X _ {-} ] e _ {n - 1} - X _ {-} X _ {+} e _ {n - 1} \\ & \qquad = H e _ {n - 1} - X _ {-} (\lambda - n + 2) e _ {n - 2} \\ & \qquad = (\lambda - 2 n + 2 + (n - 1) (\lambda - n + 2)) e _ {n - 1} \\ & \qquad = n (\lambda - n + 1) e _ {n - 1}. \end{array}
\]

推论. \(\mathbf{E}\) 的由 \(e\) 生成的子模是由诸 \(e_n\) 生成的向量子空间。

这由公式 (2) 得出。

整数 \(n \geq 0\) 中使 \(e_{n} \neq 0\) 者构成 N 中一个区间，而相应的元素 \(e_{n}\) 在 k 上构成由 e 生成的子模的基（事实上，它们线性无关，因为它们是互异权的非零元素）。称此基为与本原元 e 相伴的基。

命题 2. 若由本原元 e 生成的 E 的子模 V 是有限维的，则：

\begin{enumerate}
\def\labelenumi{(\roman{enumi})}
\item
  e 的权 \(\lambda\) 是整数且等于 \(\dim V - 1\) ；
\item
  \((e_0, e_1, \ldots, e_\lambda)\) 是 V 的基，且 \(e_n = 0\) 对 \(n > \lambda\) 成立；
\item
  \(H_{\mathrm{V}}\) 的特征值是 \(\lambda, \lambda - 2, \lambda - 4, \ldots, -\lambda\) ；它们的重数均为 1；
\item
  V 的每个本原元都与 e 成比例；
\item
  模 \(\mathrm{V}\) 的交换子可约化为纯量；特别地，\(\mathrm{V}\) 绝对单。
\end{enumerate}

设 \(m\) 是使 \(e_m \neq 0\) 的最大整数。则 \(0 = X_{+}e_{m+1} = (\lambda - m)e_m\) ，故 \(\lambda = m\) ；由于 \((e_0, e_1, \ldots, e_m)\) 是 \(V\) 的基，这就证明了 (i) 与 (ii)。断言 (iii) 由等式 \(He_n = (\lambda - 2n)e_n\) 得出。我们有 \(X_{+}e_n \neq 0\) 对 \(1 \leq n \leq m\) 成立，故得 (iv)。设 \(c\) 是模 \(V\) 的交换子的元素。则 \(Hc(e) = cH(e) = \lambda c(e)\) ，故存在 \(\mu \in k\) 使得 \(c(e) = \mu e\) ；于是

\[
c X _ {-} ^ {q} e = X _ {-} ^ {q} c e = \mu X _ {-} ^ {q} e
\]

对一切 \(q \geq 0\) 成立，故 \(c = \mu.1\) ，这就证明了 (v)。

推论. 设 E 是有限维 \(\mathfrak{sl}(2,k)\) -模。

\begin{enumerate}
\def\labelenumi{(\roman{enumi})}
\item
  自同态 \(H_{E}\) 可对角化，其特征值是有理整数。
\item
  对任何 \(p \in Z\) ，设 \(E_{p}\) 是 \(H_{E}\) 的对应于特征值 p 的特征子空间。设 i 是整数 \(\geq 0\) 。映射 \(X_{-E}^{i}|E_{p}:E_{p} \to E_{p-2i}\) 在 \(i \leq p\) 时单，在 i = p 时双，在 \(i \geq p\) 时满。映射 \(X_{+E}^{i}|E_{-p}:E_{-p} \to E_{-p+2i}\) 在 \(i \leq p\) 时单，在 i = p 时双，在 \(i \geq p\) 时满。
\item
  \(\mathrm{E}\) 的长度等于 \(\dim \operatorname{Ker} X_{+\mathrm{E}}\) ，也等于 \(\dim \operatorname{Ker} X_{-\mathrm{E}}\) 。
\item
  设 \(E'\) （相应地 \(E''\) ）是对 p 为偶（相应地奇）的诸 \(E_{p}\) 之和。则 \(E'\) （相应地 \(E''\) ）是 E 的奇（相应地偶）维单子模之和；且 \(E = E' \oplus E''\) 。\(E'\) 的长度是 \(\dim E_{0}\) ，\(E''\) 的是 \(\dim E_{1}\) 。
\end{enumerate}

\[
\text {(v)} \operatorname{Ker} X _ {+ \mathrm{E}} \cap \operatorname{Im} X _ {+ \mathrm{E}} \subset \sum_ {p > 0} \mathrm{E} _ {p} \text {and} \operatorname{Ker} X _ {- \mathrm{E}} \cap \operatorname{Im} X _ {- \mathrm{E}} \subset \sum_ {p <   0} \mathrm{E} _ {p}.
\]

若 E 单，则 E 由一个本原元生成（引理 4），只需应用命题 1 与 2。一般情形由此推出，因为每个有限维 \(\mathfrak{sl}(2,k)\) -模都半单。

\subsection{单模 V(m)}

设 \((u,v)\) 是 \(k^2\) 的典则基。对 \(\mathfrak{sl}(2,k)\) 的恒同表示，

\[
X _ {+} u = 0 \quad H u = u \quad X _ {-} u = - v
\]

\[
X _ {+} v = u \quad H v = - v \quad X _ {-} v = 0.
\]

考察对称代数 \(\mathbf{S}(k^2)\) （\(k^2\) 的）（《代数学》第三章 §6 第 1 节 定义 1）。\(\mathfrak{sl}(2,k)\) 的元素唯一地开拓为 \(\mathbf{S}(k^2)\) 的导子，从而给 \(\mathbf{S}(k^2)\) 以 \(\mathfrak{sl}(2,k)\) -模结构（第一章 §3 第 2 节）。设 \(\mathrm{V}(m)\) 是 \(\mathbf{S}(k^2)\) 中次数为 \(m\) 的齐次元素之集。则 \(\mathrm{V}(m)\) 是 \(\mathfrak{sl}(2,k)\) -子模（\(\mathbf{S}(k^2)\) 的），维数为 \(m + 1\) ，即 \(m\) 次对称幂（\(\mathrm{V}(1) = k^2\) 的）（第三章附录）。若 \(m,n\) 是满足 \(0 \leq n \leq m\) 的整数，置

\[
e _ {n} ^ {(m)} = \binom {m} {n} u ^ {m - n} v ^ {n} \in \mathrm{V} (m).
\]

命题 3. 对任何整数 \(m \geq 0\) ，\(V(m)\) 是绝对单 \(\mathfrak{sl}(2, k)\) -模。在此模中，\(e_0^{(m)} = u^m\) 是权为 \(m\) 的本原元。

我们有 \(X_{+}u^{m}=0\) 与 \(Hu^{m}=mu^{m}\) ，故 \(u^{m}\) 是权 m 的本原元。\(V(m)\) 的由 \(u^{m}\) 生成的子模维数为 \(m+1\) （命题 2 (i)），故等于 \(V(m)\) 。由命题 2 (v)，\(V(m)\) 绝对单。

定理 1. 每个单 \(\mathfrak{sl}(2,k)\) -模（维数 \(n\) 有限）都同构于 \(\mathrm{V}(n - 1)\) 。每个有限维 \(\mathfrak{sl}(2,k)\) -模都是同构于诸模 \(\mathrm{V}(m)\) 的子模的直和。

这由引理 4 与命题 1、2、3 得出。

注. 1) \(\mathfrak{sl}(2,k)\) 的伴随表示在 \(\mathfrak{sl}(2,k)\) 上定义了单 \(\mathfrak{sl}(2,k)\) -模结构。此模同构于 V(2)，同构映射把 \(u^2\) 映到 \(X_+\) ，把 \(2uv\) 映到 \(-H\) ，把 \(v^2\) 映到 \(X_-\) 。

\begin{enumerate}
\def\labelenumi{\arabic{enumi})}
\setcounter{enumi}{1}
\tightlist
\item
  对 \(n \geq 0\) 且 m \textgreater{} n，
\end{enumerate}

\[
X _ {-} e _ {n} ^ {(m)} = - (m - n) \binom{m}{n} u ^ {m - n - 1} v ^ {n + 1} = - (n + 1) e _ {n + 1} ^ {(m)}.
\]

因此，\((e_0^{(m)}, e_1^{(m)}, \ldots, e_m^{(m)})\) 是 \(V(m)\) 的与本原元 \(e_0^{(m)}\) 相伴的基。

\begin{enumerate}
\def\labelenumi{\arabic{enumi})}
\setcounter{enumi}{2}
\tightlist
\item
  设 \(\Phi\) 是 \(V(m)\) 上的双线性型，使得
\end{enumerate}

\[
\begin{array}{c} \varPhi (e _ {n} ^ {(m)}, e _ {n ^ {\prime}} ^ {(m)}) = 0 \text {if} n + n ^ {\prime} \neq m \\ \varPhi (e _ {n} ^ {(m)}, e _ {m - n} ^ {(m)}) = (- 1) ^ {n} \binom {m} {n}. \end{array}
\]

若 \(x = au + bv\) 且 \(y = cu + dv\) ，则 \(\Phi(x^{m}, y^{m}) = (ad - bc)^{m}\) 。现在容易验证 \(\Phi\) 是不变的，且 \(\Phi\) 当 m 为偶时对称，当 m 为奇时交错。

命题 4. 设 E 是有限维 \(\mathfrak{sl}(2,k)\) -模，\(m\) 是整数 \(\geq 0\) ，\(\mathrm{P}_m\) 是权为 \(m\) 的本原元之集。设 L 是从 \(\mathfrak{sl}(2,k)\) -模 \(\mathrm{V}(m)\) 到 \(\mathfrak{sl}(2,k)\) -模 E 的同态组成的向量空间。映射 \(f \mapsto f(u^m)\) （从 L 到 E 的）是线性的、单射的，其像为 \(\mathrm{P}_m \cup \{0\}\) 。

此映射显然线性，且它是单射，因为 \(u^{m}\) 生成 \(\mathfrak{sl}(2,k)\) -模 \(\mathrm{V}(m)\) 。若 \(f \in L\) ，

\[
X _ {+} (f (u ^ {m})) = f (X _ {+} u ^ {m}) = 0, \quad H (f (u ^ {m})) = f (H u ^ {m}) = m f (u ^ {m})
\]

故 \(f(u^{m}) \in \mathrm{P}_{m} \cup \{0\}\) 。设 \(e \in \mathrm{P}_{m}\) ，\(V\) 是 \(E\) 的由 \(e\) 生成的子模。由命题 1，存在从模 \(V(m)\) 到模 \(V\) 的同构，把 \(u^{m}\) 映到 \(e\) 。于是 \(L(u^{m}) = P_{m} \cup \{0\}\) 。

推论. E 的型为 V(m) 的同构型分支的长度为

\(\dim (\mathrm{P}_m\cup \{0\} .\)

\subsection{群 SL(2, k) 的线性表示}

回顾（《代数学》第三章 §8 第 9 节）我们用 \(\mathbf{SL}(2,k)\) 表示行列式等于 1 的、系数在 \(k\) 中的 2 阶方阵组成的群。若 \(x\in \mathfrak{sl}(2,k)\) 幂零，则 \(x^{2} = 0\) （《代数学》第七章 §5 命题 5 的推论 3），且 \(e^x = 1 + x\in \mathbf{SL}(2,k)\) 。若 E 是有限维向量空间而 \(\rho\) 是 \(\mathfrak{sl}(2,k)\) 在 E 上的线性表示，则 \(\rho (x)\) 幂零，于是 \(e^{\rho (x)}\) 有定义（第一章 §6 第 3 节）。

定义 2. 设 E 是有限维向量空间，\(\rho\) （相应地 \(\pi\) ）是 \(\mathfrak{sl}(2,k)\) （相应地 \(\mathbf{SL}(2,k)\) ）在 E 上的线性表示。称 \(\rho\) 与 \(\pi\) 相容，若对 \(\mathfrak{sl}(2,k)\) 的每个幂零元素 x，\(\pi(e^{x}) = e^{\rho(x)}\) 。

换言之，\(\rho\) 与 \(\pi\) 相容，若对 \(\mathfrak{sl}(2,k)\) 的每个幂零元素 x，\(\rho\) 在 kx 上的限制与 \(\pi\) 在群 \(1 + kx\) 上的限制相容（第七章 §3 第 1 节）。

若 \(\rho\) 与 \(\pi\) 相容，则 \(\rho\) 与 \(\pi\) 的对偶表示、m 次张量幂以及 m 次对称幂也分别相容（第七章 §5 第 4 节 引理 1 (i) 与 (ii)）。对 \(\rho\) 与 \(\pi\) 在 \(\rho\) 与 \(\pi\) 下稳定的向量子空间上诱导的表示，同样成立（同前引）。

特别地，表示 \(\rho_{m}\) （\(\mathfrak{sl}(2,k)\) 在 \(V(m)\) 上的，见第 3 节）与 \(\pi_{m}\) （恒同表示 \(\pi_{1}\) （\(\mathbf{SL}(2,k)\) 的）的 m 次对称幂）相容。如上置 \(e_{n}^{(m)} = \binom{m}{n} u^{m-n} v^{n}\) ，我们有

\[
\pi_ {m} (s) e _ {n} ^ {(m)} = \binom {m} {n} (s u) ^ {m - n} (s v) ^ {n}\tag{3}
\]

对 \(s \in \mathbf{SL}(2, k)\) 与 \(0 \leq n \leq m\) 成立。

定理 2. 设 \(\rho\) 是 \(\mathfrak{sl}(2,k)\) 在有限维向量空间 E 上的线性表示。

\begin{enumerate}
\def\labelenumi{(\roman{enumi})}
\item
  存在唯一的线性表示 \(\pi\) （\(\mathbf{SL}(2,k)\) 在 E 上的），它与 \(\rho\) 相容。
\item
  向量子空间 \(\mathrm{F}\) （含于 \(\mathrm{E}\) 的）在 \(\pi\) 下稳定，当且仅当它在 \(\rho\) 下稳定。
\item
  设 \(x \in \mathrm{E}\) 。则 \(\pi(s)x = x\) 对一切 \(s \in \mathbf{SL}(2,k)\) 成立，当且仅当 \(x\) 在 \(\rho\) 下不变（即 \(\rho(a)x = 0\) 对一切 \(a \in \mathfrak{sl}(2,k)\) 成立）。
\end{enumerate}

\(\pi\) 的存在性由前述及定理 1 得出。另一方面，我们知道群 \(\mathbf{SL}(2,k)\) 由形如以下的元素生成：

\[
e ^ {t X _ {+}} = \left( \begin{array}{c c} 1 & t \\ 0 & 1 \end{array} \right) \quad e ^ {- t X _ {-}} = \left( \begin{array}{c c} 1 & 0 \\ t & 1 \end{array} \right)
\]

其中 \(t \in k\) （《代数学》第三章 §8 第 9 节 命题 17）。这就证明了 \(\pi\) 的唯一性。

断言 (ii) 与 (iii) 由前述以及第七章 §3 第 1 节 引理 1 (i) 得出。证毕。

因此每个有限维 \(\mathfrak{sl}(2,k)\) -模都有唯一的 \(\mathbf{SL}(2,k)\) -模结构，称之为与其 \(\mathfrak{sl}(2,k)\) -模结构相伴的结构。

注. 当 k 是 R、C 或完备非离散超度量域时，\(\mathfrak{sl}(2,k)\) 是 \(\mathbf{SL}(2,k)\) 的李代数。设 \(\rho\) 与 \(\pi\) 如定理 2。同态 \(\pi\) 是从 \(\mathbf{SL}(2,k)\) 到 \(\mathbf{GL}(\mathrm{E})\) 的李群同态：当 \(\mathrm{E} = \mathrm{V}(m)\) 时这是显然的，一般情形由定理 1 得出。由第七章 §3 第 1 节，\(\rho(X_{+}) = \mathrm{L}(\pi)(X_{+})\) ，\(\rho(X_{-}) = \mathrm{L}(\pi)(X_{-})\) 。故 \(\rho = \mathrm{L}(\pi)\) （其逆见习题 18）。

命题 5. 设 E、F 是有限维 \(\mathfrak{sl}(2,k)\) -模，并设 \(f \in \operatorname{Hom}_{k}(\operatorname{E},\operatorname{F})\) 。下列条件等价：

\begin{enumerate}
\def\labelenumi{(\roman{enumi})}
\item
  \(f\) 是 \(\mathfrak{sl}(2,k)\) -模的同态；
\item
  \(f\) 是 \(\mathbf{SL}(2,k)\) -模的同态。
\end{enumerate}

条件 (i) 意味着 \(f\) 是 \(\mathfrak{sl}(2,k)\) -模 \(\mathrm{Hom}_k(\mathrm{E},\mathrm{F})\) 的不变元，条件 (ii) 意味着 \(f\) 是 \(\mathbf{SL}(2,k)\) -模 \(\mathrm{Hom}_k(\mathrm{E},\mathrm{F})\) 的不变元。由于这两种模结构依第七章 §5 第 4 节 引理 1 (iii) 是相伴的，本命题由定理 2 (iii) 得出。

定义 3. 群 \(\mathbf{SL}(2,k)\) 的伴随表示是 \(\mathbf{SL}(2,k)\) 在 \(\mathfrak{sl}(2,k)\) 上由下式定义的线性表示 Ad：

\[
\operatorname{Ad} (s). a = s a s ^ {- 1}
\]

对一切 \(a \in \mathfrak{sl}(2, k)\) 与一切 \(s \in \mathbf{SL}(2, k)\) 成立。

当 k 是 R、C 或完备非离散超度量域时，我们重新得到第三章 §3 第 12 节 定义 7（参见该处命题 49）。

由第七章 §5 第 4 节 引理 1 (i) 与 (ii)，\(\mathfrak{sl}(2,k)\) 与 \(\mathbf{SL}(2,k)\) 的伴随表示相容。由第七章 §3 第 1 节 注 2，\(\operatorname{Ad}(\mathbf{SL}(2,k)) = \operatorname{Aut}_e(\mathfrak{sl}(2,k))\) 。

\subsection{SL(2, k) 的一些元素}

对任何 \(t\in k^{*}\) ，置

\[
\begin{array}{l} \theta (t) = e ^ {t X _ {+}} e ^ {t ^ {- 1} X _ {-}} e ^ {t X _ {+}} \\ \qquad = \left( \begin{array}{c c} 1 & t \\ 0 & 1 \end{array} \right) \left( \begin{array}{c c} 1 & 0 \\ - t ^ {- 1} & 1 \end{array} \right) \left( \begin{array}{c c} 1 & t \\ 0 & 1 \end{array} \right) \\ \qquad = \left( \begin{array}{c c} 0 & t \\ - t ^ {- 1} & 0 \end{array} \right) \\ \qquad = e ^ {t ^ {- 1} X _ {-}} e ^ {t X _ {+}} e ^ {t ^ {- 1} X _ {-}}. \end{array}
\]

用第 3 节的记号，

\[
\theta (t) u = - t ^ {- 1} v \quad \theta (t) v = t u
\]

于是

\[
\theta (t) e _ {n} ^ {(m)} = (- 1) ^ {m - n} t ^ {2 n - m} e _ {m - n} ^ {(m)}.\tag{4}
\]

于是元素 \(\theta(t)^2 = \left( \begin{array}{cc} -1 & 0 \\ 0 & -1 \end{array} \right)\) 以 \((-1)^m\) 在 \(\mathrm{V}(m)\) 上作用。若 \(\mathrm{E}\) 是奇维单 \(\mathfrak{sl}(2,k)\) -模，则 \(\theta(t)_{\mathrm{E}}\) 是向量空间 \(\mathrm{E}\) 的对合自同构。特别地，取 \(\mathrm{E}\) 为伴随表示：

\[
\theta (t) _ {\mathrm{E}} X _ {+} = t ^ {- 2} X _ {-} \quad \theta (t) _ {\mathrm{E}} X _ {-} = t ^ {2} X _ {+} \quad \theta (t) _ {\mathrm{E}} H = - H\tag{5}
\]

于是 \(\theta(1)_{\mathrm{E}} = \theta(-1)_{\mathrm{E}}\) 是 \(\mathfrak{sl}(2,k)\) 的典则对合。

对任何 \(t\in k^{*}\) ，置

\[
h (t) = \left( \begin{array}{c c} t & 0 \\ 0 & t ^ {- 1} \end{array} \right) = \theta (t) \theta (- 1).
\]

则 \(h(t)u = tu\) ，\(h(t)v = t^{-1}v\) ，于是

\[
h (t) e _ {n} ^ {(m)} = t ^ {m - 2 n} e _ {n} ^ {(m)}.\tag{6}
\]

命题 6. 设 E 是有限维 \(\mathfrak{sl}(2,k)\) -模，且 \(t\in k^{*}\) 。设 \(\mathrm{E}_p\) 是 E 中权为 \(p\) 的元素之集。

\begin{enumerate}
\def\labelenumi{(\roman{enumi})}
\item
  \(\theta(t)_{\mathrm{E}}|\mathrm{E}_{p}\) 是从 \(\mathrm{E}_p\) 到 \(\mathrm{E}_{-p}\) 的双射。
\item
  \(h(t)_{\mathrm{E}}|\mathrm{E}_p\) 是比为 \(t^p\) 的位似（在 \(\mathrm{E}_p\) 上的）。
\end{enumerate}

若 \(\mathrm{E} = \mathrm{V}(n)\) ，本命题由公式 (4) 与 (6) 得出。一般情形由定理 1 得出。

推论. 设 \(E = E' \oplus E''\) 是命题 2 的推论中定义的 E 的分解。元素 \(\begin{pmatrix}-1 & 0 \\ 0 & -1\end{pmatrix}\) （\(\mathbf{SL}(2, k)\) 的）在 \(E'\) 上以 +1 作用，在 \(E''\) 上以 -1 作用。

这由应用于 \(t = -1\) 的 (ii) 得出。

\section{分裂半单李代数的根系}

\subsection{分裂半单李代数}

定义 1. 设 \(\mathfrak{g}\) 是半单李代数。嘉当子代数 \(\mathfrak{h}\) （\(\mathfrak{g}\) 的）称为分裂的，若对一切 \(x \in \mathfrak{h}\) ，\(\operatorname{ad}_{\mathfrak{g}} x\) 可三角化。半单李代数称为可分裂的，若它有一个分裂嘉当子代数。分裂半单李代数是指一个偶 \((\mathfrak{g}, \mathfrak{h})\) ，其中 \(\mathfrak{g}\) 是半单李代数，而 \(\mathfrak{h}\) 是 \(\mathfrak{g}\) 的分裂嘉当子代数。

注. 1) 设 \(\mathfrak{g}\) 是半单李代数，\(\mathfrak{h}\) 是 \(\mathfrak{g}\) 的嘉当子代数。对一切 \(x \in \mathfrak{h}\) ，\(\operatorname{ad}_{\mathfrak{g}} x\) 半单（第七章 §2 第 4 节 定理 2）。于是，说 \(\mathfrak{h}\) 分裂就意味着 \(\operatorname{ad}_{\mathfrak{g}} x\) 对一切 \(x \in \mathfrak{h}\) 可对角化。

\begin{enumerate}
\def\labelenumi{\arabic{enumi})}
\setcounter{enumi}{1}
\item
  若 \(k\) 代数闭，则每个半单李代数 \(\mathfrak{g}\) 都可分裂，且 \(\mathfrak{g}\) 的每个嘉当子代数都分裂。当 \(k\) 不是代数闭时，存在不可分裂的半单李代数（习题 2 a)）；此外，若 \(\mathfrak{g}\) 可分裂，\(\mathfrak{g}\) 也可能有不分裂的嘉当子代数（习题 2 b)）。
\item
  设 \(\mathfrak{g}\) 是半单李代数，\(\mathfrak{h}\) 是 \(\mathfrak{g}\) 的嘉当子代数，\(\rho\) 是 \(\mathfrak{g}\) 的有限维单射表示且使 \(\rho(\mathfrak{h})\) 可对角化。则 \(\operatorname{ad}_{\mathfrak{g}}x\) 对一切 \(x \in \mathfrak{h}\) 可对角化（第七章 §2 第 1 节 例 2），故 \(\mathfrak{h}\) 分裂。
\item
  我们将看到（§3 第 3 节 命题 10 的推论），若 \(\mathfrak{h}\) 与 \(\mathfrak{h}'\) 是 \(\mathfrak{g}\) 的分裂嘉当子代数，则存在 \(\mathfrak{g}\) 的初等自同构把 \(\mathfrak{h}\) 变为 \(\mathfrak{h}'\) 。
\item
  设 \(\mathfrak{g}\) 是约化李代数。则 \(\mathfrak{g} = \mathfrak{c} \times \mathfrak{s}\) ，其中 \(\mathfrak{c}\) 是 \(\mathfrak{g}\) 的中心，而 \(\mathfrak{s} = \mathcal{D}\mathfrak{g}\) 半单。\(\mathfrak{g}\) 的嘉当子代数正是形如 \(\mathfrak{h} = \mathfrak{c} \times \mathfrak{h}'\) 的子代数，其中 \(\mathfrak{h}'\) 是 \(\mathfrak{s}\) 的嘉当子代数（第七章 §2 第 1 节 命题 2）。此时称 \(\mathfrak{h}\) 分裂，若 \(\mathfrak{h}'\) 相对于 \(\mathfrak{s}\) 分裂。这以一种明显的方式引出可分裂或分裂约化代数的定义。
\end{enumerate}